\documentclass[11pt]{article}
\usepackage[T1]{fontenc}
\usepackage[utf8]{inputenc}
\usepackage{lmodern}
\usepackage{amsmath,amssymb,amsthm,mathtools}
\usepackage[margin=1.05in]{geometry}
\usepackage{microtype,booktabs,array,enumitem}
\usepackage{newunicodechar}
\newunicodechar{∃}{\ensuremath{\exists}}
\newunicodechar{∀}{\ensuremath{\forall}}
\newunicodechar{ℝ}{\ensuremath{\mathbb R}}
\newunicodechar{ℕ}{\ensuremath{\mathbb N}}
\newunicodechar{∧}{\ensuremath{\wedge}}
\newunicodechar{≤}{\ensuremath{\le}}
\newunicodechar{ε}{\ensuremath{\varepsilon}}
\usepackage[dvipsnames]{xcolor}
\usepackage[colorlinks=true,linkcolor=blue!45!black,citecolor=blue!45!black,urlcolor=blue!45!black]{hyperref}
\usepackage[capitalise,nameinlink]{cleveref}
\hypersetup{pdftitle={A linear bound for the Paulsen problem},pdfauthor={Astra, Claude, Hugo Eberhard, Kaarel H\"anni}}

\newtheorem{theorem}{Theorem}[section]
\newtheorem{lemma}[theorem]{Lemma}
\newtheorem{proposition}[theorem]{Proposition}
\newtheorem{corollary}[theorem]{Corollary}
\theoremstyle{definition}
\newtheorem{definition}[theorem]{Definition}
\theoremstyle{remark}
\newtheorem{remark}[theorem]{Remark}

\newcommand{\R}{\mathbb R}
\newcommand{\E}{\mathbb E}
\newcommand{\Prob}{\mathbb P}
\newcommand{\one}{\mathbf 1}
\newcommand{\F}{\mathrm F}
\newcommand{\op}{\mathrm{op}}
\newcommand{\eps}{\varepsilon}
\newcommand{\ip}[2]{\langle #1,#2\rangle}
\newcommand{\norm}[1]{\lVert #1\rVert}
\newcommand{\fro}[1]{\lVert #1\rVert_{\F}}
\newcommand{\opn}[1]{\lVert #1\rVert_{\op}}
\newcommand{\Lap}{\mathcal L}
\DeclareMathOperator{\diag}{diag}
\DeclareMathOperator{\Diag}{Diag}
\DeclareMathOperator{\tr}{tr}
\DeclareMathOperator{\Cov}{Cov}
\DeclareMathOperator{\Var}{Var}
\DeclareMathOperator{\rank}{rank}
\DeclareMathOperator{\osc}{osc}
\DeclareMathOperator{\im}{im}
\numberwithin{equation}{section}
\setlength{\emergencystretch}{2em}
\allowdisplaybreaks[1]

\title{A linear bound for the Paulsen problem}
\author{Astra\thanks{AI system (gpt-astra).}
\and Claude\thanks{AI system, Anthropic.}
\and Hugo Eberhard
\and Kaarel H\"anni}
\date{October 2026}

\begin{document}
\maketitle
\vspace{-1.5em}
\begin{abstract}
We prove that every real $\eps$-nearly equal-norm Parseval frame of $n$ vectors
in $\R^d$ is within squared Frobenius distance $C\eps d$ of an equal-norm
Parseval frame, where $C$ is an absolute constant; this is optimal up to the
value of $C$. Equivalently, every rank-$d$ orthogonal projection on $\R^n$ whose
diagonal entries deviate from $d/n$ by at most $\beta$ is within squared
Frobenius distance $Cn\beta$ of a rank-$d$ orthogonal projection with constant
diagonal. The proof has two parts. A deterministic \emph{static balancing
theorem}, proved by minimising a log-determinant potential over a box, provides
a row scaling that corrects every diagonal error that is small compared with the
inverse of the \emph{barrier constant}: a half-set hitting time of the random
walk whose jump rates are the squared entries $P_{ij}^2$ of the projection. Two random \emph{seeds} then make the barrier
constant small at linear cost: a filtered Gaussian tangent perturbation with a
deterministic drift when $d\gtrsim\log n$, and a row-wise tangent perturbation
with row renormalisation when $n\gtrsim d^2$. Bounded $d$ is handled by the
Hamilton--Moitra argument. The main theorem has been formalised in Lean~4 with
Mathlib.
\end{abstract}

\tableofcontents

\section{The result and the plan}\label{sec:intro}

\subsection{Statement}
A frame of $n$ vectors in $\R^d$ is a matrix $X\in\R^{n\times d}$ whose $i$th
row is $x_i^T$. Distances are \emph{squared} Frobenius distances
$\fro{X-Y}^2=\sum_i\norm{x_i-y_i}^2$, with no normalisation. Throughout,
\[
 a=d/n .
\]
$X$ is \emph{Parseval} if $X^TX=I_d$ and \emph{equal-norm} if
$\norm{x_i}^2=a$ for all $i$; an equal-norm Parseval frame is an \emph{ENP frame}.
$X$ is an \emph{$\eps$-nearly ENP frame} if
\begin{equation}\label{eq:nearly}
 (1-\eps)I_d\preceq X^TX\preceq(1+\eps)I_d,
 \qquad (1-\eps)a\le\norm{x_i}^2\le(1+\eps)a\quad(1\le i\le n).
\end{equation}

\begin{theorem}[Linear Paulsen bound]\label{thm:main}
There is an absolute constant $C$ such that for all $1\le d\le n$, all
$0<\eps<1$ and every real $\eps$-nearly ENP frame $X\in\R^{n\times d}$ there is
an ENP frame $W\in\R^{n\times d}$ with $\fro{X-W}^2\le C\eps d$.
\end{theorem}

The bound is optimal up to the constant (\cref{rem:optimal}). Kwok, Lau, Lee and
Ramachandran~\cite{KLLR} proved the first bound independent of $n$,
$O(\eps d^{13/2})$; Hamilton and Moitra~\cite{HM} improved it to
$O(\eps d^2)$; Lau and Ramachandran announced the linear bound~\cite{LR}, via a
different route. \Cref{app:formal} describes a Lean~4 formalisation of
\cref{thm:main,thm:projection}.

For a Parseval $U$ write $P=UU^T$ for the orthogonal projection onto its column
space; then $P_{ii}=\norm{u_i}^2$. The problem has an equivalent form in which
the error is additive.

\begin{theorem}[Projection form]\label{thm:projection}
There is an absolute constant $C$ such that for every $n\ge1$ and every
rank-$d$ orthogonal projection $P$ on $\R^n$, with
$\beta=\max_i|P_{ii}-d/n|$, there is a rank-$d$ orthogonal projection $Q$ with
$Q_{ii}=d/n$ for all $i$ and $\fro{P-Q}^2\le Cn\beta$.
\end{theorem}

This reformulation and the complement symmetry are due to Cahill and
Casazza~\cite{CC}: both $\beta$ and $\fro{P-Q}$ are unchanged when $(P,Q)$ is
replaced by $(I-P,I-Q)$, so in \cref{thm:projection} one may assume $d\le n/2$.
In \cref{sec:assembly}, \cref{thm:projection} is derived from an equal-norm
formulation (\cref{prop:equalrow}) and \cref{thm:main} from
\cref{thm:projection}.

\begin{remark}[Optimality]\label{rem:optimal}
Let $d=2d_1$, $n=2n_1$, and $1\le k\le n_1/4$ with $n_1-k\ge d_1$; put
$\eps=4k/n$. Let $X=X_1\oplus X_2$, where $X_1$ is an ENP frame of $n_1-k$
vectors in $\R^{d_1}$ and $X_2$ one of $n_1+k$ vectors in $\R^{d_1}$
(\cref{cor:exist}). Then $X$ is Parseval, and its squared row norms
$a(1\mp2k/n)^{-1}$ make it an $\eps$-nearly ENP frame. Let $W$ be ENP, and let $E$
and $A$ be the principal blocks of $XX^T$ and $WW^T$ on the first $n_1-k$ rows.
Then $E$ is a projection of trace $d_1$ and $A$ a positive contraction of trace
$(n_1-k)a$, so $\Delta=E-A$ has $\tr\Delta=ka=\eps d/4$. Since
$\tr(E\Delta)=\tr\Delta+\tr((I-E)A)\ge\tr\Delta$,
\[
 \tr(A-A^2)=-\tr\Delta+2\tr(E\Delta)-\tr\Delta^2\ge\tr\Delta-\fro\Delta^2 .
\]
As $WW^T$ is a projection, its off-diagonal block has squared norm $\tr(A-A^2)$,
while that of $XX^T$ vanishes; hence
$\fro{XX^T-WW^T}^2\ge\fro\Delta^2+2(\tr\Delta-\fro\Delta^2)\ge2\tr\Delta-\fro{XX^T-WW^T}^2$,
i.e.\ $\fro{XX^T-WW^T}^2\ge\tr\Delta$, and by \cref{lem:align}
$\fro{X-W}^2\ge\frac12\tr\Delta=\eps d/8$.
\end{remark}

\paragraph{Conventions.} $c,C$ denote positive absolute constants whose value
may change between occurrences; constants with subscripts or names are fixed
once chosen. $\opn{\cdot}$ is the Euclidean operator norm, $\one$ the all-ones
vector, $\Diag(x)$ the diagonal matrix with diagonal $x$, $\diag(M)$ the diagonal
of $M$ as a vector, and $\Pi_{\mathcal S}$ the orthogonal projection onto a
subspace $\mathcal S$. $A\preceq B$ is the Loewner order. All Gaussian vectors are
real and centred.

\subsection{The plan}\label{subsec:plan}

\paragraph{The engine: scaling.}
Let $U$ be Parseval with projection $P$ and diagonal $p=\diag P$. For
$s\in\R^n$ the matrix $\Diag(e^s)U$ has a column space with projection $P(s)$,
and $s\mapsto\diag P(s)$ is the gradient of the convex function
$\frac12\log\det(U^T\Diag(e^{2s})U)$, whose Hessian at $s=0$ is $2L_P$, where
\[
 L_P=\Diag(p)-P\circ P
\]
is the graph Laplacian with edge weights $P_{ij}^2$ (\cref{lem:energy}). To
first order, changing the diagonal by $\delta\perp\one$ therefore requires
$s\approx\frac12L_P^{+}\delta$, at a cost of about
$s^TL_Ps\approx\frac14\delta^TL_P^{+}\delta$. If $L_P$ has spectral gap
$\lambda$ and $|\delta_i|\le\eps a$, the cost is at most $n\eps^2a^2/\lambda$.
For a ``random-looking'' $P$ one has $\lambda\approx a$ and the cost is
$\eps^2 d$, but $L_P$ can be arbitrarily badly connected: for an exact direct
sum the components cannot exchange diagonal mass at all.

\paragraph{Perturb, then scale.}
Perturb $U$ by a random horizontal tangent of size $t=\sqrt{K_0\eps}$, with
$K_0$ a large constant, at cost $O(t^2d)=O(\eps d)$. Generic entries of the new
projection then have size $t\sqrt{a/n}$, so most edges acquire weight
$\gtrsim at^2/n$, and the random walk with rates $P_{ij}^2$ reaches any half of
the vertices in time $O(1/(at^2))$: the \emph{barrier constant}
(\cref{def:barrier,lem:core}) is $O(1/(at^2))$. The static balancing theorem
(\cref{thm:balancing}) then removes exactly every diagonal error $\beta$ with
$\beta\cdot O(1/(at^2))\le\frac12$, at cost $\frac{15}4n\beta=O(t^2d)$. The
$\ell_\infty$ (rather than spectral) formulation is essential: the scaling must
stay in a box $\norm{s}_\infty\le1$, and an $\ell_2$ argument loses a factor
$\sqrt n$.

\paragraph{The difficulty.}
The perturbed frame must therefore have diagonal error at most $\theta at^2$ in
$\ell_\infty$, for a small fixed $\theta$. The initial error $\eps a=at^2/K_0$
is acceptable. But a random tangent of size $t$ moves each diagonal entry
\emph{linearly} by about $t\sqrt{a/n}$, which exceeds $at^2$ as soon as
$\eps\lesssim1/d$. Two different constructions deal with this.

\paragraph{Seed 1: many rows ($n\ge Bd^2$, \cref{sec:manyrow}).}
Here the input is an equal-norm frame $X$ with spectral error
$\eta=\opn{X^TX-I}$, and $t^2$ is a large multiple of $\eta$. Perturb each row
\emph{orthogonally to itself} and renormalise. Row norms stay exact, but the
frame operator moves. Its first-order change is removed by conditioning the
Gaussian on a subspace of codimension at most $d^2$, and its second-order
fluctuation is an average over $n$ rows, of size $t^2d/\sqrt n$. The spectral
error of the perturbed frame is therefore $\eta+O(t^2d/\sqrt n+t^2d^2/n+t^3)$,
which is at most $\theta t^2$ once $t^2/\eta$ and $B$ are large and $t$ is
small; whitening turns a spectral error $\delta$ into a diagonal error at most
$2a\delta$. The cubic term is $O(t^3)$, not the termwise $O(t^3d)$, because the
first-order constraint allows the normalisation weights to be centred
(\cref{lem:remainder}); this is why the seed works for every $\eta\le c_*$,
with no relation required between $\eta$ and $d$.

\paragraph{Seed 2: moderate rows ($d\ge A\log 2n$, \cref{sec:moderate}).}
Perturb $U$ in the horizontal tangent space of the Grassmannian, so that the
frame stays exactly Parseval, with a Gaussian whose covariance is filtered:
$C_\rho=\rho(I-\Omega)(\rho I+\Omega)^{-1}$, where $\Omega$ is the normalised
Gram operator of the linear diagonal map $\mathcal A$ and $\rho=D^2t^2$ for a
large constant $D$. The filter makes the linear diagonal change
$O(t\sqrt{\rho a\log n/n})=O(Dat^2\sqrt{\log n/d})$, negligible when
$d\gg\log n$, while leaving the weak (connecting) directions untouched. The
price is a second-order bias $t^2b_*$, which is not small in $\ell_\infty$. A
commutator identity shows that $b_*=\mathcal Am_*$ for an explicit
deterministic tangent direction $m_*$ with $\fro{m_*}^2\le2a/\rho$
(\cref{lem:commutator,lem:mean}). We add the drift $\frac t2m_*$ to the noise:
its first-order effect cancels the bias exactly, and its other effects are
$O(at^2/D)$. All remaining errors are at most $\theta at^2$ by concentration,
which is where $d\ge A\log 2n$ is used. Only a bounded, deterministic set of
rows needs a spectral argument, and second moments suffice there.

\paragraph{Assembly (\cref{sec:assembly}).}
After complementation and normalisation it suffices to treat equal-norm frames
with $2d\le n$ and spectral error $\eta$. If $d<d_0$ (a large absolute
constant), the Hamilton--Moitra bound $\eta d(d-1)\le d_0\eta d$ suffices
(\cref{sec:bounded}). If $d\ge d_0$ and $n\ge Bd^2$, Seed~1 applies. If
$d\ge d_0$ and $n<Bd^2$, then $A\log 2n\le A\log(2Bd^2)\le d$ and Seed~2
applies.

\begin{center}
\small
\begin{tabular}{@{}lll@{}}
\toprule
regime & construction & where\\
\midrule
$d<d_0$ & radial isotropic position, ordered coordinates (Hamilton--Moitra) & \cref{sec:bounded}\\
$d\ge d_0,\ n\ge Bd^2$ & row-tangent Gaussian, row normalisation, whitening, balancing & \cref{sec:manyrow}\\
$d\ge d_0,\ 2d\le n<Bd^2$ & filtered horizontal Gaussian with drift, balancing & \cref{sec:moderate}\\
$d>n/2$ & complement & \cref{sec:assembly}\\
\bottomrule
\end{tabular}
\end{center}

\section{The scaling toolbox}\label{sec:toolbox}

Throughout this section $U\in\R^{n\times d}$ is Parseval, $P=UU^T$,
$p=\diag P$, and $L_P=\Diag(p)-P\circ P$. For nonnegative symmetric weights
$w_{ij}$ on $[n]$ the weighted Laplacian $L$ acts by
$(Lx)_i=\sum_jw_{ij}(x_i-x_j)$; diagonal weights do not contribute. Unless
another Laplacian is specified, $L=L_P$. We write $\osc(x)=\max_ix_i-\min_ix_i$.

\subsection{Projections, distances, and the squared-Gram Laplacian}

\begin{lemma}[Frames versus projections]\label{lem:align}
Let $U,W\in\R^{n\times d}$ be Parseval with projections $P,Q$, and let
$\sigma_1,\dots,\sigma_d\in[0,1]$ be the singular values of $U^TW$. Then
\[
 \min_{R\in O(d)}\fro{U-WR}^2=2\sum_k(1-\sigma_k)\le
 \fro{P-Q}^2=2\sum_k(1-\sigma_k^2)=2\fro{(I-P)W}^2\le2\fro{U-W}^2 .
\]
Right multiplication by $R\in O(d)$ preserves Parsevalness and every row norm.
If $X\in\R^{n\times d}$ has $\delta=\opn{X^TX-I}<1$, its polar factor
$\tilde X=X(X^TX)^{-1/2}$ is Parseval, $\fro{X-\tilde X}^2\le d\delta^2$, and
$\norm{x_i}^2/(1+\delta)\le\norm{\tilde x_i}^2\le\norm{x_i}^2/(1-\delta)$. If
moreover $\delta\le\frac12$, then $|\norm{\tilde x_i}^2-\norm{x_i}^2|\le2\delta\norm{x_i}^2$
and row $i$ of $\tilde X\tilde X^T-XX^T$ has squared norm at most
$6\delta^2\norm{x_i}^2$.
\end{lemma}

\begin{proof}
$\fro{U-WR}^2=2d-2\tr(U^TWR)$, maximised over $R$ by a polar factor of $U^TW$ (any orthogonal $R$ with $U^TWR\succeq0$; it need not be unique),
giving $2d-2\sum_k\sigma_k$. Next $\fro{P-Q}^2=2d-2\tr(PQ)=2d-2\fro{U^TW}^2$
and $\fro{(I-P)W}^2=d-\fro{U^TW}^2$; also $1-\sigma\le1-\sigma^2$ on $[0,1]$,
and $(I-P)W=(I-P)(W-U)$ gives $\fro{(I-P)W}\le\fro{U-W}$. For the polar factor, if $\varsigma_k$ are the singular values of $X$
then $\fro{X-\tilde X}^2=\sum_k(\varsigma_k-1)^2\le\sum_k(\varsigma_k^2-1)^2$, and
$\norm{\tilde x_i}^2=x_i^T(X^TX)^{-1}x_i$. With $E=(X^TX)^{-1}-I$, so
$\opn E\le\delta/(1-\delta)\le2\delta$, row $i$ of $\tilde X\tilde X^T-XX^T=XEX^T$
is $XEx_i$, of squared norm at most $(1+\delta)4\delta^2\norm{x_i}^2$.
\end{proof}

\begin{lemma}[Energy identity]\label{lem:energy}
$L$ is the Laplacian of the weights $P_{ij}^2$ $(i\ne j)$, $L\succeq0$,
$L\one=0$, $L_{I-P}=L_P$, and for $x\in\R^n$, $X=\Diag(x)$,
\[
 x^TLx=\tfrac12\sum_{i,j}P_{ij}^2(x_i-x_j)^2=\fro{(I-P)XU}^2
 =\tfrac12\fro{[X,P]}^2 .
\]
\end{lemma}

\begin{proof}
$\sum_jP_{ij}^2=(P^2)_{ii}=p_i$, so the rows of $\Diag(p)-P\circ P$ sum to zero
and its off-diagonal entries are $-P_{ij}^2$. Next
$\fro{(I-P)XU}^2=\tr(XPX(I-P))=\sum_ix_i^2p_i-\sum_{i,j}x_ix_jP_{ij}^2$, and
$[X,P]_{ij}=(x_i-x_j)P_{ij}$. For $I-P$ the off-diagonal weights are again
$P_{ij}^2$.
\end{proof}

\begin{lemma}[Scaled projections]\label{lem:scaled}
For $w\in\R^n_{>0}$ put $D=\Diag(w)$, $M=(U^TD^2U)^{-1/2}$ and $W=DUM$. Then
$W$ is Parseval with column space $D\im U$, and
\[
 \fro{WW^T-P}^2=2\fro{(I-P)W}^2\le\frac{2\,w^TLw}{\min_iw_i^2} .
\]
\end{lemma}

\begin{proof}
$W^TW=MU^TD^2UM=I$ and $U^TD^2U\succeq(\min_iw_i)^2I$, so
$\opn M\le1/\min_iw_i$. By \cref{lem:align,lem:energy},
$\fro{(I-P)W}^2\le\opn M^2\fro{(I-P)DU}^2=\opn M^2\,w^TLw$.
\end{proof}

\subsection{Diagonal scaling}

For $s\in\R^n$ let $D_s=\Diag(e^{s})$, let $P(s)$ be the projection onto
$D_s\im U$, $r(s)=\diag P(s)$, and $F(s)=\frac12\log\det(U^TD_s^2U)$.

\begin{lemma}\label{lem:potential}
$F$ is smooth and convex, $\nabla F(s)=r(s)$, $\sum_ir_i(s)=d$, and
$F(s+c\one)=F(s)+cd$.
\end{lemma}

\begin{proof}
With $G=U^TD_s^2U=\sum_ie^{2s_i}u_iu_i^T$ one has
$\partial_iF=e^{2s_i}u_i^TG^{-1}u_i=P(s)_{ii}$, because $P(s)=D_sUG^{-1}U^TD_s$;
differentiating again gives $\nabla^2F=2L_{P(s)}\succeq0$. The trace of $P(s)$ is
$d$, and $G$ scales by $e^{2c}$ under $s\mapsto s+c\one$.
\end{proof}

\begin{lemma}[One-sided scaling inequalities]\label{lem:scaling-ineq}
For $z\in\R^n_{>0}$ let $r$ be the diagonal of the projection onto
$\Diag(\sqrt z)\im U$. Then, coordinatewise,
\[
 Lz\le z\circ(r-p),\qquad L(z^{-1})\le-z^{-1}\circ(r-p).
\]
\end{lemma}

\begin{proof}
Let $G=U^T\Diag(z)U$. Since $G+z_i^2G^{-1}-2z_iI=(G^{1/2}-z_iG^{-1/2})^2\succeq0$,
evaluating at $u_i$ and using $u_i^TGu_i=\sum_jP_{ij}^2z_j=p_iz_i-(Lz)_i$ and
$z_i^2u_i^TG^{-1}u_i=z_ir_i$ gives the first inequality. Let
$V\in\R^{n\times(n-d)}$ be an isometry onto $(\im U)^\perp$. The subspaces
$\Diag(\sqrt z)\im U$ and $\Diag(z^{-1/2})\im V$ are orthogonal complements, so
the projection onto the latter has diagonal $\one-r$, while $VV^T=I-P$ has
diagonal $\one-p$ and Laplacian $L$ (\cref{lem:energy}). The first inequality
for $(V,z^{-1})$ is the second. (If $d=n$ then $L=0$ and $r=p=\one$.)
\end{proof}

\subsection{The barrier constant and static balancing}

\begin{definition}[Barrier constant]\label{def:barrier}
Let $L$ be a weighted Laplacian on $[n]$. For $S\subseteq[n]$, an
\emph{$S$-barrier} is a vector $\psi\ge0$ with $(L\psi)_i\ge1$ for every
$i\notin S$. The \emph{barrier constant} $H(L)\in(0,\infty]$ is the least $H$ such
that every $S$ with $|S|\ge n/2$ has an $S$-barrier with $\norm\psi_\infty\le H$
($H(L)=\infty$ if some such $S$ has no barrier).
\end{definition}

For the continuous-time random walk that jumps from $i$ to $j$ at rate
$w_{ij}$, the expected hitting time $h_S(x)=\E_x\tau_S$ of $S$ vanishes on $S$
and satisfies $(Lh_S)_i=1$ off $S$. If $\psi$ is any $S$-barrier and $\beta'>1$,
the maximum principle argument in the proof of \cref{lem:median}, applied to
$h_S-\beta'\psi$, gives $h_S\le\beta'\psi$. So $h_S$ is the minimal $S$-barrier,
and $H(L_P)$ is the half-set hitting time of the walk with rates $P_{ij}^2$: the
largest expected time, over all starting points and all $S$ with $|S|\ge n/2$,
needed to reach $S$. If $H(L)<\infty$ then the graph is connected: for a
component $J$ with $|J|\le n/2$, summing $(L\psi)_i$ over $i\in J$ gives $0$, so
$S=[n]\setminus J$ has no barrier. The barrier constant is within a factor $2$ of
the $\ell_\infty$ Poisson constant of $L$ (\cref{rem:poisson}). A small barrier
constant controls scalings in $\ell_\infty$:

\begin{lemma}[Median barrier]\label{lem:median}
Let $L$ be a weighted Laplacian on $[n]$ with $H(L)\le H<\infty$, let $\beta\ge0$ with $\beta H<1$, let
$z\in\R^n_{>0}$, and let $m$ be a median of $z$ (so that $\{z\le m\}$ and
$\{z\ge m\}$ both have at least $n/2$ elements). If
\[
 (Lz)_i\le\beta z_i\ \text{ when }z_i>m,\qquad
 (Lz^{-1})_i\le\beta z_i^{-1}\ \text{ when }z_i<m ,
\]
then $(1-\beta H)\max_iz_i\le m\le(1-\beta H)^{-1}\min_iz_i$.
\end{lemma}

\begin{proof}
Let $S=\{z\le m\}$, let $\psi$ be an $S$-barrier with $\norm\psi_\infty\le H$, let
$M=\max z$ and $\beta'>\beta$. The vector $v=z-m\one-\beta'M\psi$ is $\le0$ on
$S$, and for $i\notin S$, $(Lv)_i\le\beta z_i-\beta'M<0$. At a maximiser $i$ of
$v$, $(Lv)_i=\sum_jw_{ij}(v_i-v_j)\ge0$; so every maximiser lies in $S$, and
$v\le0$. Hence $M\le m+\beta'MH$; let $\beta'\downarrow\beta$. Apply the same
argument to $z^{-1}$, whose median is $m^{-1}$, with $S=\{z\ge m\}$.
\end{proof}

\begin{theorem}[Static balancing]\label{thm:balancing}
Suppose $H(L_P)\le H<\infty$, and $q\in\R^n$ satisfies $\sum_iq_i=d$ and
$\beta H\le\frac12$, where $\beta=\norm{q-p}_\infty$. Then
there is a Parseval $W$ with $\diag(WW^T)=q$ and
\[
 \fro{U-W}^2\le\fro{UU^T-WW^T}^2\le\tfrac{15}4\norm{q-p}_1 .
\]
$W$ is obtained from $U$ by a row scaling with ratio of scales at most $2$,
whitening, and an orthogonal right factor.
\end{theorem}

\begin{proof}
Let $s$ minimise $\Phi(s)=F(s)-\ip qs$ over the cube $[0,1]^n$ (compactness),
and write $r=r(s)$. Since $\partial_i\Phi=r_i-q_i$ (\cref{lem:potential}),
coordinatewise optimality gives
\begin{equation}\label{eq:kkt}
 r_i\le q_i\ \text{ if }s_i>0,\qquad r_i\ge q_i\ \text{ if }s_i<1 .
\end{equation}
Let $z=e^{2s}$ with median $m$, so $1\le\min z\le m\le\max z\le e^{2}$. If
$z_i>m$ then $s_i>0$, so by \cref{lem:scaling-ineq} and~\eqref{eq:kkt}
$(Lz)_i\le z_i(r_i-p_i)\le z_i(q_i-p_i)\le\beta z_i$; if $z_i<m$ then $s_i<1$
and likewise $(Lz^{-1})_i\le\beta z_i^{-1}$. By \cref{lem:median},
$\max z/\min z\le(1-\beta H)^{-2}\le4$, i.e.\ $\osc(s)\le\log2<1$. Hence
$s$ does not meet both faces of the cube: either all $s_i>0$, and then
$r\le q$, or some $s_j=0$, all $s_i<1$, and then $r\ge q$. Since
$\sum r_i=d=\sum q_i$, in both cases $r=q$.

For the distance let $w=e^{s}$, $f=q-p$, and let $c$ be the midpoint of the
range of $z^2$. \Cref{lem:scaling-ineq} with $r=q$ gives $Lz\le z\circ f$, so,
using $\sum_if_i=0$,
\[
 z^TLz\le\sum_iz_i^2f_i=\sum_i(z_i^2-c)f_i\le\tfrac12\osc(z^2)\norm f_1 .
\]
As $(z_i-z_j)^2\ge4\min w^2\,(w_i-w_j)^2$, \cref{lem:energy,lem:scaled} give
\[
 \fro{P(s)-P}^2\le\frac{2\,w^TLw}{\min w^2}\le\frac{z^TLz}{2\min z^2}
 \le\frac{\osc(z^2)}{4\min z^2}\norm f_1\le\frac{4^2-1}4\norm f_1 .
\]
Finally apply \cref{lem:align} to the whitened frame $DU(U^TD^2U)^{-1/2}$,
$D=\Diag(w)$, whose projection is $P(s)$.
\end{proof}

\begin{remark}
The hypothesis forces $q_i\ge p_i/2>0$: for $n\ge2$, the barrier for
$S=[n]\setminus\{i\}$ gives $1\le(L\psi)_i\le p_i(1-p_i)H$, so
$\beta\le1/(2H)\le p_i/2$. The proof uses no existence theorem for scalings,
no flow, and no genericity or connectivity argument: the box supplies a
minimiser, the median barrier keeps it away from one face, and the trace
identity $\sum_ir_i=d$ does the rest. With
$\beta H<1$ and a box of side $T>-\log(1-\beta H)$ the same proof gives scale
ratio $(1-\beta H)^{-1}$ and cost $\frac14((1-\beta H)^{-4}-1)\norm{q-p}_1$.
\end{remark}

\subsection{Barriers from a dense core}

\begin{lemma}[Dense core with exceptional vertices]\label{lem:core}
Let $L$ be the Laplacian of nonnegative symmetric weights $w_{ij}$ on $[n]$,
and let $[n]=\mathcal G\sqcup\mathcal B$. Suppose:
\begin{enumerate}[label=(\alph*)]
\item there is $\kappa>0$ with $\sum_{j\in S}w_{ij}\ge\kappa$ for every
$i\in\mathcal G$ and every $S\subseteq[n]\setminus\{i\}$ with $|S|\ge n/2$;
\item there are $y\in\R^{\mathcal B}$, $y\ge0$, and $R,F_{\mathcal B}\ge0$ with
$\norm y_\infty\le R$, $(L_{\mathcal B\mathcal B}y)_i\ge1$ for $i\in\mathcal B$,
and $\sum_{j\in\mathcal B}w_{ij}y_j\le F_{\mathcal B}$ for $i\in\mathcal G$.
(If $\mathcal B=\varnothing$ take $R=F_{\mathcal B}=0$.)
\end{enumerate}
Then $H(L)\le(1+F_{\mathcal B})/\kappa+R$. Moreover, with $b=|\mathcal B|$:
\begin{enumerate}[label=(\roman*)]
\item (a) holds with $\kappa=\vartheta\gamma$ if $\vartheta\in(0,\frac12]$,
$\gamma>0$, and every $i\in\mathcal G$ has at least $(\frac12+\vartheta)n$ indices
$j\ne i$ with $w_{ij}\ge\gamma/n$;
\item (b) holds with $R=\sqrt b/\lambda$, $F_{\mathcal B}=b$, if
$x^TLx\ge\lambda\norm x^2$ for every $x$ supported on $\mathcal B$ ($\lambda>0$);
\item for $w_{ij}=P_{ij}^2$, (b) holds with $R=\frac8{3\alpha}$,
$F_{\mathcal B}=\frac2{3\alpha}\max_ip_i$ if $\alpha>0$, $p_i\ge\alpha/2$ on
$\mathcal B$ and $\sum_{i\in\mathcal B}p_i\le\frac14$.
\end{enumerate}
\end{lemma}

\begin{proof}
Let $|S|\ge n/2$, put $\lambda_0=(1+F_{\mathcal B})/\kappa$, extend $y$ by $0$ to
$\mathcal G$, and let $\psi=(\lambda_0\one+y)\circ\one_{S^c}$, so
$0\le\psi\le\lambda_0+R$ and $\psi_j\le\lambda_0+y_j$ for all $j$. For
$i\in\mathcal G\setminus S$, $\psi_i=\lambda_0$ and $\psi_j=0$ on $S$, so
\[
 (L\psi)_i\ge\lambda_0\sum_{j\in S}w_{ij}-\sum_{j\in\mathcal B}w_{ij}y_j\ge
 \lambda_0\kappa-F_{\mathcal B}=1 .
\]
For $i\in\mathcal B\setminus S$, $\psi_i=\lambda_0+y_i$, so
$(L\psi)_i\ge\sum_jw_{ij}(y_i-y_j)=(L_{\mathcal B\mathcal B}y)_i\ge1$. Thus $\psi$
is an $S$-barrier.

(i) If $N_i$ is the set of such $j$ and $S\subseteq[n]\setminus\{i\}$, then
$|N_i\cap S|\ge|N_i|+|S|-(n-1)>\vartheta n$.
(ii) $\mathsf L=L_{\mathcal B\mathcal B}\succeq\lambda I$; let $y=\mathsf L^{-1}\one_{\mathcal B}$. If
$y$ had a negative minimum at $i$, then $(\mathsf Ly)_i=y_i\sum_{j\notin\mathcal
B}w_{ij}+\sum_{j\in\mathcal B}w_{ij}(y_i-y_j)\le0$, a contradiction; so $y\ge0$,
and $\norm y_\infty\le\norm y_2\le\sqrt b/\lambda$. Since
$(\mathsf L\one_{\mathcal B})_j=\sum_{i\in\mathcal G}w_{ij}$ for $j\in\mathcal B$, for each
$i_0\in\mathcal G$
\[
 \sum_{j\in\mathcal B}w_{i_0j}y_j\le\sum_{i\in\mathcal G}\sum_{j\in\mathcal B}w_{ij}y_j
 =\one_{\mathcal B}^T\mathsf Ly=b .
\]
(iii) $P_{ij}^2\le p_ip_j$ gives, for $i\in\mathcal B$,
$(L_{\mathcal B\mathcal B}\one_{\mathcal B})_i=p_i-\sum_{j\in\mathcal B}P_{ij}^2\ge
p_i(1-\sum_{j\in\mathcal B}p_j)\ge\frac{3\alpha}8$, and for $i\in\mathcal G$,
$\sum_{j\in\mathcal B}P_{ij}^2\le p_i/4$. Take $y=\frac8{3\alpha}\one_{\mathcal B}$.
\end{proof}

\begin{remark}[Poisson constants]\label{rem:poisson}
The $\ell_\infty$ Poisson constant $K$ of $L$ is the least $K$ such that every
$f\perp\one$ has a solution of $Lg=f$ with $\norm g_\infty\le K\norm f_\infty$. Then $K\le H\le2K$. Indeed, for
$|S|\ge n/2$ the forcing $f=\one-\frac n{|S|}\one_S$ has $\norm f_\infty\le1$,
and $g+K\one$ is an $S$-barrier; conversely, if $H<\infty$ the graph is connected,
so $Lg=f$ is solvable for every $f\perp\one$, and if $\norm f_\infty\le1$ and $m$ is
a median of $g$, the argument of \cref{lem:median}
applied to $g-m\one-\beta'\psi$ ($\beta'>1$) gives $g\le m+H$, and symmetrically
$g\ge m-H$. Barriers are more convenient than Poisson solutions because they
can be written down explicitly, as in \cref{lem:core}.
\end{remark}

\subsection{From a seed to an exact frame}

Both constructions produce the following object.

\begin{definition}[Seed]\label{def:seed}
Let $\Theta\ge2$ and $t>0$. A \emph{$(\Theta,t)$-seed} for $X\in\R^{n\times d}$
is a Parseval $V\in\R^{n\times d}$ with
\[
 \fro{V-X}^2\le\Theta t^2d,\qquad
 H(L_{VV^T})\le\frac{\Theta}{at^2},\qquad
 \max_i\bigl|(VV^T)_{ii}-a\bigr|\le\frac{at^2}{2\Theta}.
\]
\end{definition}

\begin{corollary}\label{cor:seed}
If $X$ has a $(\Theta,t)$-seed, there is an ENP frame $W$ with
$\fro{X-W}^2\le3\Theta t^2d$.
\end{corollary}

\begin{proof}
Apply \cref{thm:balancing} to $V$ with $q=a\one$: $\beta H\le\frac12$, and
$\fro{V-W}^2\le\frac{15}4n\cdot\frac{at^2}{2\Theta}\le t^2d$. Then
$\fro{X-W}^2\le2\fro{X-V}^2+2\fro{V-W}^2\le(2\Theta+2)t^2d\le3\Theta t^2d$.
\end{proof}

\section{Bounded rank: the Hamilton--Moitra estimate}\label{sec:bounded}

This section reproduces the radial-isotropic and ordered-coordinate argument of
Hamilton and Moitra~\cite{HM}, with a direct compactness proof of the existence
of the radial scaling. It is used for $d$ below a large
absolute constant, and it also shows that ENP frames exist.

\begin{lemma}\label{lem:HM}
Let $1\le d\le n$, let $X\in\R^{n\times d}$ have $\norm{x_i}^2=a$ for all $i$,
and let $\eta=\opn{X^TX-I}$. There is an ENP frame $W$ with
$\fro{X-W}^2\le\eta d(d-1)$.
\end{lemma}

\begin{proof}
If $d=1$ then $X^TX=na=1$ and $W=X$ works. Let $d\ge2$.

\emph{Step 1: full spark.} Suppose first that every $d\times d$ minor
$\mu_S=\det(X_S)^2$ ($S\subset[n]$, $|S|=d$) is nonzero. Consider
$F(s)=\frac12\log\det(X^T\Diag(e^{2s})X)-a\sum_is_i$. By Cauchy--Binet,
\[
 F(s)=\tfrac12\log\sum_{|S|=d}\mu_S\,e^{T(s)_S},\qquad
 T(s)_S=2\Bigl(\sum_{i\in S}s_i-a\sum_{i}s_i\Bigr),
\]
where $T:\R^n\to\R^{\binom nd}$ is linear and $\sum_ST(s)_S=0$ because each $i$
lies in a fraction $d/n=a$ of the $d$-sets. The function
$\mathcal E(y)=\sum_S\mu_Se^{y_S}$ restricted to the subspace $\im T$ has the
nonempty sublevel set $\{\mathcal E\le\mathcal E(0)\}$, on which
$y_S\le\log(\mathcal E(0)/\mu_S)$ for every $S$ and hence, by $\sum_Sy_S=0$,
also $y_S$ is bounded below. This set is compact, so $\mathcal E|_{\im T}$
attains its minimum at some $T(s^*)$, and $s^*$ minimises $F$.

Let $D=\Diag(e^{s^*})$, $M=(X^TD^2X)^{-1/2}$ and $W=DXM$. The computation in
the proof of \cref{lem:potential} (which does not use $X^TX=I$) gives
$\partial_iF=\norm{w_i}^2-a$, so $W$ is ENP, and since
$w_i=e^{s^*_i}Mx_i$ has norm $\sqrt a$,
\[
 w_i=\sqrt a\,\frac{Mx_i}{\norm{Mx_i}}.
\]

\emph{Step 2: ordered coordinates.} Replacing $X,W$ by $XO,WO$ for a suitable
$O\in O(d)$ changes neither the hypotheses nor the distance, and makes
$M=\Diag(\lambda_1,\dots,\lambda_d)$ with $0<\lambda_1\le\dots\le\lambda_d$.
Then $w_{ij}=\lambda_jx_{ij}/\sqrt{c_i}$ with
$c_i=a^{-1}\sum_\ell\lambda_\ell^2x_{i\ell}^2$, so $w_{ij}$ has the sign of
$x_{ij}$ and $w_{ij}^2-x_{ij}^2=x_{ij}^2(\lambda_j^2/c_i-1)$ is nonpositive and
then nonnegative as $j$ increases. As $\sum_j(x_{ij}^2-w_{ij}^2)=a-a=0$, the
prefix sums $\Delta_i(k)=\sum_{j\le k}(x_{ij}^2-w_{ij}^2)$ are nonnegative,
vanish at $k=d$, and
\[
 \sum_j|x_{ij}^2-w_{ij}^2|=2\max_k\Delta_i(k)\le2\sum_{k=1}^{d-1}\Delta_i(k).
\]
Since $W^TW=I$ and $(X^TX)_{jj}\le1+\eta$,
$\sum_i\Delta_i(k)=\sum_{j\le k}((X^TX)_{jj}-1)\le k\eta$. Equal signs give
$(x_{ij}-w_{ij})^2\le|x_{ij}^2-w_{ij}^2|$, and therefore
\[
 \fro{X-W}^2\le2\eta\sum_{k=1}^{d-1}k=\eta d(d-1).
\]

\emph{Step 3: general $X$.} Let $G$ be a Vandermonde matrix with rows
$(1,\theta_i,\dots,\theta_i^{d-1})$ for distinct $\theta_i$; all its maximal
minors are nonzero. Each maximal minor of $X+hG$ is a polynomial in $h$ with
nonzero leading coefficient, so for all but finitely many $h$ the matrix $X+hG$
is full spark; choose such $h_k\to0$ (so small that no row of $X+h_kG$
vanishes) and normalise each row of $X+h_kG$ to length $\sqrt a$ (a positive diagonal left factor, which preserves full spark).
The resulting $X_k\to X$ have ENP frames $W_k$ with
$\fro{X_k-W_k}^2\le\eta_kd(d-1)$, $\eta_k\to\eta$. The set of ENP frames is
compact; a limit point $W$ of $(W_k)$ is ENP and satisfies the bound.
\end{proof}

\begin{corollary}\label{cor:exist}
For all $1\le d\le n$ an ENP frame exists. Consequently every equal-norm
$X\in\R^{n\times d}$ is within $\fro{X-W}^2\le2\fro X^2+2\fro W^2=4d$ of an ENP
frame $W$.
\end{corollary}

\begin{proof}
Apply \cref{lem:HM} to $x_i=\sqrt a\,e_1$ for all $i$.
\end{proof}

\section{The many-row seed}\label{sec:manyrow}

\begin{theorem}\label{thm:manyrow}
There are absolute constants $B,c_*,C_*>0$ such that the following holds. Let
$2\le d$, $n\ge Bd^2$, let $X\in\R^{n\times d}$ have $\norm{x_i}^2=a=d/n$ for all
$i$, and suppose $\eta:=\opn{X^TX-I}$ satisfies $0<\eta\le c_*$. Then there is
an ENP frame $W$ with $\fro{X-W}^2\le C_*\eta d$.
\end{theorem}

Throughout, $S=X^TX$, so
$\opn{S-I}\le\eta\le c_*\le\frac12$ (the choice below gives $c_*\le\frac1{64}$), and
$e_i=x_i/\sqrt a$.

\subsection{The noise}

In the Frobenius space $\R^{n\times d}$ consider the subspaces
\[
 \mathcal E=\{Z:\ip{x_i}{z_i}=0\ \forall i\},\qquad
 \mathcal F=\{Z\in\mathcal E:\ X^TZ+Z^TX=0\},
\]
and let $m=\dim\mathcal E-\dim\mathcal F\le\dim\mathrm{Sym}_d\le d^2\le n$. For a
standard Gaussian $g\in\R^{n\times d}$ put
\[
 Z_0=n^{-1/2}\Pi_{\mathcal E}g,\qquad Z=n^{-1/2}\Pi_{\mathcal F}g,\qquad
 Z_\perp=Z_0-Z=n^{-1/2}\Pi_{\mathcal E\ominus\mathcal F}g .
\]
$Z$ and $Z_\perp$ are independent. The rows $z_{0,i}=n^{-1/2}(I-e_ie_i^T)g_i$
of $Z_0$ are independent. Since $\Cov(Z)\preceq\Cov(Z_0)$, each row $z_i$ of
$Z$ is Gaussian with covariance $\Sigma_i\preceq n^{-1}(I-e_ie_i^T)$, and
$\Cov(\mathrm{vec}\,Z)\preceq I/n$. Moreover
\begin{equation}\label{eq:mr-noise}
 \E\fro Z^2=\frac{n(d-1)-m}{n}\le d,\qquad
 \E\fro{Z-Z_0}^2=\frac mn\le\frac{d^2}n .
\end{equation}

\begin{lemma}[Row moments]\label{lem:rowmoments}
Let $r_i=\norm{z_i}^2/a$, $\mu_i=\E r_i$ and $\delta_i=1-\mu_i$. Then
$1/d\le\delta_i\le1$, $\sum_ia\delta_i=1+m/n\le2$, and
\[
 \mathcal M:=a\sum_i\E\bigl[(r_i-1)^2(1+r_i)^2\bigr]\le C_{\mathcal M},
\]
with $C_{\mathcal M}$ absolute.
\end{lemma}

\begin{proof}
$\tilde\Sigma_i=\Sigma_i/a\preceq d^{-1}I$ has trace $\mu_i\le(d-1)/d$, so
$\tr\tilde\Sigma_i^2\le1/d$ and $\tr\tilde\Sigma_i^4\le d^{-3}$. For
$\xi_i=r_i-\mu_i$ the Gaussian moment formulas give $\E\xi_i^2=2\tr\tilde\Sigma_i^2\le
2/d$ and $\E\xi_i^4=12(\tr\tilde\Sigma_i^2)^2+48\tr\tilde\Sigma_i^4\le60/d^2$.
Next $\sum_ia\delta_i=d-\E\fro Z^2=1+m/n$ by~\eqref{eq:mr-noise}. Finally
$(r_i-1)^2\le2\xi_i^2+2\delta_i^2$ and $(1+r_i)^2\le8+2\xi_i^2$, so
\[
 \E(r_i-1)^2(1+r_i)^2\le16\E\xi_i^2+4\E\xi_i^4+16\delta_i^2+4\delta_i^2\E\xi_i^2
 \le\frac{C}d+C\delta_i^2 ,
\]
and $a\sum_i(C/d+C\delta_i^2)\le C+Ca\sum_i\delta_i\le3C$.
\end{proof}

\subsection{Row normalisation and the exact second-order identity}

For $0<t\le1$ set
\[
 v_i=\frac{x_i+tz_i}{\sqrt{1+t^2r_i}},\qquad
 v_{0,i}=\frac{x_i+tz_{0,i}}{\sqrt{1+t^2r_{0,i}}},\qquad
 r_{0,i}=\norm{z_{0,i}}^2/a .
\]
Since $z_i\perp x_i$, $\norm{v_i}^2=\norm{v_{0,i}}^2=a$ exactly. The map
$y\mapsto\sqrt a\,y/\norm y$ is the metric projection onto the closed ball of
radius $\sqrt a$ on $\{\norm y\ge\sqrt a\}$, hence $1$-Lipschitz there. Thus
\begin{equation}\label{eq:mr-contract}
 \fro{V-X}\le t\fro Z,\quad \fro{V-V_0}\le t\fro{Z-Z_0},\quad
 \opn V\le\opn{X+tZ},\quad\opn{V_0}\le\opn{X+tZ_0},
\end{equation}
the last two because $V=\Diag(\lambda)(X+tZ)$ with $0<\lambda_i\le1$, and similarly
for $V_0$.

Expanding $1/(1+t^2r_i)=1-t^2c_i$ with
\[
 c_i=\frac{r_i}{1+t^2r_i},
\]
and using $X^TZ+Z^TX=0$, one finds the exact identity
\begin{equation}\label{eq:mr-identity}
 V^TV=S+t^2Q(Z)+R_3+R_4,\qquad Q(Z)=Z^TZ-\sum_ir_ix_ix_i^T,
\end{equation}
\[
 R_3=-t^3\sum_ic_i(x_iz_i^T+z_ix_i^T),\qquad
 R_4=-t^4\sum_ic_i(z_iz_i^T-r_ix_ix_i^T).
\]

\begin{lemma}[Mean and fluctuation of $Q$]\label{lem:Q}
$\opn{\E Q(Z)-(I-S)}\le m/n$ and $\E\fro{Q(Z)-\E Q(Z)}^2\le8d^2/n$.
\end{lemma}

\begin{proof}
For $Z_0$: $\E Z_0^TZ_0=n^{-1}\sum_i(I-e_ie_i^T)=I-S/d$ and $\E r_{0,i}=(d-1)/d$,
so $\E Q(Z_0)=I-S/d-(1-1/d)S=I-S$. As $Z_0=Z+Z_\perp$ with independent centred
summands and $Q$ is quadratic, $\E Q(Z)=\E Q(Z_0)-\E Q(Z_\perp)$, and
$\E Q(Z_\perp)$ is the difference of the positive semidefinite matrices
$\E Z_\perp^TZ_\perp$ and $\sum_i\E(\norm{z_{\perp,i}}^2/a)x_ix_i^T$, both of trace
$\E\fro{Z_\perp}^2=m/n$; so its norm is at most $m/n$.

Entry $(\alpha,\beta)$ of $Q$ is the quadratic form $\sum_iz_i^TB_{\alpha\beta,i}z_i$
with $B_{\alpha\beta,i}=\frac12(E_{\alpha\beta}+E_{\beta\alpha})-(e_i)_\alpha(e_i)_\beta I$.
For a Gaussian vector with covariance $\Sigma\preceq I/n$ and symmetric
coefficient matrix $\mathsf B$, the variance of the quadratic form is
$2\fro{\Sigma^{1/2}\mathsf B\Sigma^{1/2}}^2\le2\fro{\mathsf B}^2/n^2$. Since
$\sum_{\alpha,\beta}\fro{B_{\alpha\beta,i}}^2\le2d^2+2d\le4d^2$ for each $i$,
summing over $\alpha,\beta$ and the $n$ blocks gives the claim.
\end{proof}

\begin{lemma}[Remainder]\label{lem:remainder}
Let $\bar c=1/(1+t^2)$, $\Gamma=\Diag(c_i-\bar c)$ and
$\Lambda=\Diag(c_ir_i-\bar c)$. Then
\[
 \opn{R_3}\le2t^3\opn X\fro{\Gamma Z},\qquad
 \opn{R_4}\le t^4\bigl(\opn Z^2+\opn Z\fro{\Gamma Z}+\opn S+\opn X\fro{\Lambda X}\bigr),
\]
and $\E\fro{\Gamma Z}^2\le\mathcal M$, $\E\fro{\Lambda X}^2\le\mathcal M$.
\end{lemma}

\begin{proof}
Because $\sum_i(x_iz_i^T+z_ix_i^T)=X^TZ+Z^TX=0$, the constant part $\bar c$ of
the weights drops out of $R_3$:
$R_3=-t^3(X^T\Gamma Z+Z^T\Gamma X)$, which gives the first bound. Next
$\sum_ic_iz_iz_i^T=\bar cZ^TZ+Z^T\Gamma Z$ and $\sum_ic_ir_ix_ix_i^T=\bar cS+X^T\Lambda X$,
with $\opn{Z^T\Gamma Z}\le\opn Z\fro{\Gamma Z}$ and similarly for $\Lambda$, and
$\bar c\le1$. For the moments,
\[
 c_i-\bar c=\frac{r_i-1}{(1+t^2r_i)(1+t^2)},\qquad
 c_ir_i-\bar c=\frac{(r_i-1)(1+r_i+t^2r_i)}{(1+t^2r_i)(1+t^2)},
\]
so $|c_i-\bar c|\le|r_i-1|$ and, as $1+r+t^2r\le(1+r)(1+t^2r)$,
$|c_ir_i-\bar c|\le|r_i-1|(1+r_i)$. Hence
$\fro{\Gamma Z}^2=a\sum_i(c_i-\bar c)^2r_i\le a\sum_i(r_i-1)^2(1+r_i)^2$ and
$\fro{\Lambda X}^2=a\sum_i(c_ir_i-\bar c)^2\le a\sum_i(r_i-1)^2(1+r_i)^2$; take
expectations and use \cref{lem:rowmoments}.
\end{proof}

\begin{remark}
The termwise bound $\opn{R_3}\le2t^3\sum_ic_i\norm{x_i}\norm{z_i}\approx2t^3d$
would only allow $\eta\lesssim d^{-2}$. Centring the weights at
$\bar c$, which is allowed by the first-order constraint, replaces the $n$
terms of size $a$ by fluctuations of relative size $d^{-1/2}$, whose
$\ell_2$-sum over the rows is $O(1)$.
\end{remark}

\subsection{A dense reference graph}

\begin{lemma}\label{lem:mr-dense}
There are absolute $h,t_0\in(0,1]$ and $c_1>0$ such that for $d\ge2$ and
$t\le t_0$, with
probability at least $1-ne^{-c_1(n-1)}$, every $i$ has at least $0.9(n-1)$
indices $j\ne i$ with
\[
 |\ip{v_{0,i}}{v_{0,j}}|\ge h\,t\sqrt{a/n}.
\]
\end{lemma}

\begin{proof}
Write $z_{0,j}/\sqrt a=d^{-1/2}g'_j$ with $g'_j\sim N(0,I-e_je_j^T)$
independent. Fix $i$ and condition on row $i$; let $f=v_{0,i}/\sqrt a$, a unit
vector. For $j\ne i$,
\[
 \frac{\ip{v_{0,i}}{v_{0,j}}}{a}=\frac{\nu_j}{N_j},\qquad
 \nu_j=u_j+\tfrac t{\sqrt d}\ip f{g'_j},\quad u_j=\ip f{e_j},\quad
 N_j=\bigl(1+\tfrac{t^2}d\norm{g'_j}^2\bigr)^{1/2}\ge1,
\]
and the target inequality reads $|\nu_j|/N_j\ge ht/\sqrt d$. By Markov,
$\Prob(N_j>2)\le t^2/3$. On $\{N_j\le2\}$ failure requires
$|\nu_j|<2ht/\sqrt d$. Here $\nu_j\sim N(u_j,t^2(1-u_j^2)/d)$. If $|u_j|\le\frac12$
the variance is at least $3t^2/(4d)$ and the Gaussian density bound gives
probability at most $4h/\sqrt{3\pi/2}\le3h$. If $|u_j|>\frac12$ and $2ht\le\frac14$,
failure requires $\frac t{\sqrt d}|\ip f{g'_j}|>\frac14$, of probability at most
$16t^2/d$. Choose $h$ and $t_0$ so that $3h+t_0^2/3+8t_0^2\le\frac1{50}$ (this also gives
$2ht\le\frac14$, and $16t^2/d\le8t^2$ as $d\ge2$). The
failure indicators for $j\ne i$ are conditionally independent with
probabilities at most $\frac1{50}$, so by the exponential Markov inequality
(with parameter $1$) more than $0.1(n-1)$ failures have probability at most
$\exp(-(n-1)[0.1-\frac{e-1}{50}])\le e^{-c_1(n-1)}$. Take a union bound over
$i$.
\end{proof}

\subsection{Proof of \texorpdfstring{\cref{thm:manyrow}}{Theorem}}

\emph{Constants.} Fix $h,t_0,c_1$ from \cref{lem:mr-dense}. Put $\kappa_0=h^2/200$,
$C'=36^2\cdot20$, and
$C_{\rm mix}=125/h^2+3$, then $\Theta=\max\{42,C_{\rm mix}\}$ and
$\theta=\min\{1/(4\Theta),h/60\}$. Let $C_R=300+100\sqrt{C_{\mathcal M}}$, choose $t_1=\min\{t_0,\theta/(16C_R)\}$, and set
\[
 t^2=4\eta/\theta,\qquad c_*=\theta t_1^2/4,
\]
so that $\eta\le c_*$ gives $t\le t_1$. Finally choose $B$ so large that
$n\ge Bd^2$ implies $d^2/n\le\theta/4$, $\sqrt{160d^2/n}\le\theta/4$,
$2aC'd/\kappa_0\le\frac14$ (this bounds $\sum_{i\in B_0}2a$ for the set $B_0$
below, which has $|B_0|\le C'd/\kappa_0$), $n\ge30$, and $ne^{-c_1(n-1)}\le0.01$.

\emph{The event.} By Markov's inequality, \eqref{eq:mr-noise},
\cref{lem:Q,lem:remainder}, the net bound (\cref{lem:gauss}(c)) for $Z$ and
$Z_0$, and \cref{lem:mr-dense}, with positive probability all of the following
hold:
\[
 \fro Z^2\le20d,\quad \fro{Z-Z_0}^2\le\frac{20d^2}n,\quad
 \fro{Q-\E Q}^2\le\frac{160d^2}n,\quad
 \fro{\Gamma Z}^2\le20\mathcal M,\quad \fro{\Lambda X}^2\le20\mathcal M,
\]
$\opn Z,\opn{Z_0}\le16$, and the conclusion of \cref{lem:mr-dense}: the five Markov
events fail with probability at most $\frac1{20}$ each, the two net bounds
(\cref{lem:gauss}(c) with $\sigma^2=1/n$, $u=16$) with probability at most
$2\cdot5^{n+d}e^{-8n}$ each, and \cref{lem:mr-dense} with probability at most
$0.01$. Fix such a sample.

\emph{Spectral error.} By~\eqref{eq:mr-identity},
$V^TV-I=(1-t^2)(S-I)+t^2(\E Q-(I-S))+t^2(Q-\E Q)+R_3+R_4$. Since
$\opn X\le\sqrt2$, $\opn S\le\frac32$ and $t\le1$, \cref{lem:remainder} and the
event give
$\opn{R_3}+\opn{R_4}\le t^3\bigl(2\sqrt{40\mathcal M}+256+16\sqrt{20\mathcal M}
+\tfrac32+\sqrt{40\mathcal M}\bigr)\le C_Rt^3$, and therefore
\[
 \delta:=\opn{V^TV-I}\le\eta+t^2\frac{d^2}n+t^2\sqrt{\frac{160d^2}n}+C_Rt^3
 \le\frac\theta4t^2+\frac\theta4t^2+\frac\theta4t^2+\frac\theta{16}t^2<\theta t^2 .
\]

\emph{Graph.} By~\eqref{eq:mr-contract}, $\opn V,\opn{V_0}\le\sqrt2+16\le18$
and
$\fro{VV^T-V_0V_0^T}\le36\,t\fro{Z-Z_0}$, so
$\fro{VV^T-V_0V_0^T}^2\le C't^2d^2/n$. Let $B_0$ be the set of rows whose
contribution to this sum exceeds $\kappa_0at^2$; then $|B_0|\le C'd/\kappa_0$. A row
$i\notin B_0$ has at most $\kappa_0at^2/((h/2)^2t^2a/n)=n/50$ indices with
$|(VV^T-V_0V_0^T)_{ij}|\ge\frac h2t\sqrt{a/n}$, so by \cref{lem:mr-dense} it has at
least $0.85n$ indices $j\ne i$ with $|(VV^T)_{ij}|\ge\frac h2t\sqrt{a/n}$.

Let $U=V(V^TV)^{-1/2}$, $P=UU^T$. By \cref{lem:align} (rows of $P-VV^T$) a row loses at most
$6a\theta^2t^4/((h/4)^2t^2a/n)\le n/20$ of these indices, so each $i\notin
B_0$ has at least $0.8n$ indices $j\ne i$ with $P_{ij}^2\ge(h^2/16)at^2/n$.
Also $a/2\le P_{ii}\le2a$ and $\sum_{i\in B_0}P_{ii}\le\frac14$.
\Cref{lem:core} with $\mathcal B=B_0$, part (i) with $\vartheta=0.3$ and
$\gamma=h^2at^2/16$, and part (iii) with $\alpha=a$ (so $R=\frac8{3a}$,
$F_{\mathcal B}\le\frac43$) shows that
$H(L_P)\le\frac{7/3}{0.3\gamma}+\frac8{3a}\le C_{\rm mix}/(at^2)$.

\emph{Seed.} $\fro{U-X}^2\le2d\delta^2+2t^2\fro Z^2\le42t^2d$, the diagonal error
is at most $2a\delta\le2\theta at^2\le at^2/(2\Theta)$, and the barrier constant
is at most $\Theta/(at^2)$. So $U$ is a $(\Theta,t)$-seed, and
\cref{cor:seed} gives an ENP frame within $3\Theta t^2d=(12\Theta/\theta)\eta d$
of $X$.
\qed

\section{The moderate-row seed}\label{sec:moderate}

\begin{theorem}\label{thm:moderate}
There are absolute constants $A,C_m,\eps_0>0$ such that the following holds.
Let $U\in\R^{n\times d}$ be Parseval with $2d\le n$, $d\ge A\log(2n)$, and
$|\norm{u_i}^2-a|\le\eps a$ for all $i$, where $0<\eps\le\eps_0$. Then there is
an ENP frame $\widehat U$ with $\fro{\widehat U-U}^2\le C_m\eps d$.
\end{theorem}

Throughout this section $P=UU^T$, $p=\diag P$, $D_p=\Diag(p)$, $L=L_P$, and
$\eps_0\le\frac12$, so $\frac a2\le p_i\le\frac{3a}2\le\frac34$. For a symmetric
$M\in\R^{n\times n}$, $\Lap(M)=\sum_{i<j}M_{ij}^2(e_i-e_j)(e_i-e_j)^T$ is the
Laplacian with weights $M_{ij}^2$; thus $L=\Lap(P)$,
$x^T\Lap(M)x=\frac12\fro{[\Diag x,M]}^2$, and
\begin{equation}\label{eq:lap-basic}
 \Lap(M)\preceq2\max_i\norm{M_i}^2\,I,\qquad
 \Lap(M+M')\succeq\tfrac12\Lap(M)-\Lap(M'),
\end{equation}
by $(x_i-x_j)^2\le2x_i^2+2x_j^2$ and $(v+w)^2\ge\frac12v^2-w^2$. We write
$\mathsf K=nI-\one\one^T$ for the Laplacian of the complete graph.

\paragraph{Scales.} Put $\sigma=at^2$, where $t\asymp\sqrt\eps$ is the size of
the perturbation. The perturbation creates edges of weight $\gtrsim\sigma/n$,
the resulting barrier constant is $O(1/\sigma)$, and so all diagonal errors must
be at most a small multiple of $\sigma$. The table lists the errors and how
each is controlled ($\rho=D^2t^2$ is the filter level, $D$ a large constant).
\begin{center}
\small
\begin{tabular}{@{}lll@{}}
\toprule
diagonal error & size$/\sigma$ & reason\\
\midrule
input & $\eps/t^2$ & $t^2=K_0\eps$, $K_0$ large\\
linear term $2t\mathcal AZ$ & $6D\sqrt{\log(2n)/d}$ & filter (\cref{lem:filter}), $d\ge A\log 2n$\\
quadratic fluctuation & $\eta'$ & concentration, $d\ge A\log 2n$\\
filter bias $t^2b_*$ & $0$ & cancelled by the drift (\cref{lem:mean,lem:retraction})\\
drift cross terms & $40/D+1/(2D^2)$ & $D$ large\\
base bias $t^2b_0$ & $12\eps/d$ & \cref{lem:mean}\\
retraction remainder & $C_Et$ & $t$ small\\
\bottomrule
\end{tabular}
\end{center}

\subsection{The horizontal tangent space and the filtered covariance}

Fix an isometry $V\in\R^{n\times(n-d)}$ onto $(\im U)^\perp$, and let
$\mathcal Z=\R^{(n-d)\times d}$ with the Frobenius inner product. For
$Z\in\mathcal Z$ the matrix $VZ$ is horizontal ($U^TVZ=0$), and we put
\[
 Y_Z=VZU^T+UZ^TV^T,\qquad \norm{(Y_Z)_i}^2=\norm{Z^Tv_i}^2+\norm{Zu_i}^2\le\opn Z^2,
 \qquad \fro{Y_Z}^2=2\fro Z^2 .
\]
Define
\[
 \mathcal A:\mathcal Z\to\R^n,\quad (\mathcal AZ)_i=v_i^TZu_i=(VZU^T)_{ii},\qquad
 \mathcal A^*x=N_x:=V^T\Diag(x)U .
\]
Then $\mathcal A\mathcal A^*=L$, i.e.\ $\fro{N_x}^2=x^TLx$ (\cref{lem:energy}). Let
\[
 \bar{\mathcal A}=D_p^{-1/2}\mathcal A,\qquad
 S=\bar{\mathcal A}\bar{\mathcal A}^*=D_p^{-1/2}LD_p^{-1/2},\qquad
 \Omega=\bar{\mathcal A}^*\bar{\mathcal A} .
\]
Since $L\succeq0$ and $S=I-D_p^{-1/2}(P\circ P)D_p^{-1/2}$ with $P\circ P\succeq0$
(Schur product theorem),
\begin{equation}\label{eq:S}
 0\preceq S\preceq I,\qquad 0\preceq \Omega\preceq I,\qquad
 \tr(I-S)=\sum_i\frac{P_{ii}^2}{p_i}=d .
\end{equation}
Fix an orthonormal eigenbasis of $S$, and for each of its members $y$ with
eigenvalue $\mu>0$ put $f_y=D_p^{-1/2}y$ and
$Z_y=\mu^{-1/2}\bar{\mathcal A}^*y=\mu^{-1/2}N_{f_y}$. Since $\Omega Z_y=\mu Z_y$ and
$\ip{Z_y}{Z_{y'}}=\ip y{Sy'}/\sqrt{\mu\mu'}$, the $Z_y$ form an orthonormal
eigenbasis of $\Omega$ on $(\ker \Omega)^\perp$. Sums $\sum_{\mu>0}$ below run over these
basis members (with multiplicity). Since $p_i\ge a/2$,
\begin{equation}\label{eq:finfty}
 \norm{f_y}_\infty\le\sqrt{2/a}.
\end{equation}

\begin{definition}[Filtered covariance]
For $\rho>0$ let $C_\rho=\rho(I-\Omega)(\rho I+\Omega)^{-1}$ on $\mathcal Z$, let
$Z\sim N(0,C_\rho/n)$, and $Y=Y_Z$.
\end{definition}

\begin{lemma}\label{lem:filter}
\begin{enumerate}[label=(\alph*)]
\item $0\preceq C_\rho\preceq I$, and $C_\rho=I$ on $\ker \Omega$.
\item $\Cov(\mathcal AZ)=n^{-1}D_p^{1/2}\,\rho S(I-S)(\rho I+S)^{-1}D_p^{1/2}\preceq
(\rho/n)D_p$.
\item $I-C_\rho=\Omega+R_\rho$ with $R_\rho=\sum_{\mu>0}r_\mu\,Z_y\otimes Z_y$,
$r_\mu=\mu(1-\mu)/(\rho+\mu)\ge0$, and
\[
 \sum_{\mu>0}r_\mu\le d,\qquad \sum_{\mu>0}\frac{r_\mu}{\sqrt\mu}\le\frac d{2\sqrt\rho},
 \qquad\sum_{\mu>0}\frac{r_\mu}\mu\le\frac d\rho .
\]
\end{enumerate}
\end{lemma}

\begin{proof}
(a) The eigenvalues of $C_\rho$ are $\rho(1-\mu)/(\rho+\mu)\in[0,1]$. (b) By
singular value calculus $\bar{\mathcal A}C_\rho\bar{\mathcal A}^*=\rho S(I-S)(\rho
I+S)^{-1}$, whose eigenvalues $\rho\mu(1-\mu)/(\rho+\mu)$ are at most $\rho$. (c)
On the $\mu$-eigenspace, $1-\rho(1-\mu)/(\rho+\mu)=\mu+\mu(1-\mu)/(\rho+\mu)$. The
sums follow from $\mu/(\rho+\mu)\le1$, $\sqrt\mu/(\rho+\mu)\le1/(2\sqrt\rho)$,
$1/(\rho+\mu)\le1/\rho$, and $\sum_{\mu>0}(1-\mu)\le\tr(I-S)=d$.
\end{proof}

\begin{remark}
Only these properties of the filter $\phi(\mu)=\rho(1-\mu)/(\rho+\mu)$ are used:
$\phi(0)=1$ (so $C_\rho=I$ on $\ker \Omega$), $0\le\phi(\mu)\le1-\mu$,
$\mu\phi(\mu)\le\rho$, and the sums in (c). The hard cutoff $(1-\mu)\one[\mu<\rho]$
satisfies them with $\sum r_\mu/\sqrt\mu\le d/\sqrt\rho$, which doubles some
constants below; the rational filter keeps every object a rational function of
$S$, which is convenient for formalisation.
\end{remark}

\subsection{Commutator identities}

For $Z\in\mathcal Z$ define the quadratic diagonal and its polarisation
\[
 q_i(Z)=\norm{Z^Tv_i}^2-\norm{Zu_i}^2,\qquad
 q_i(Z,M)=\ip{Z^Tv_i}{M^Tv_i}-\ip{Zu_i}{Mu_i},
\]
so that $q(Z+M)=q(Z)+2q(Z,M)+q(M)$ and, by Cauchy--Schwarz,
\begin{equation}\label{eq:q-cs}
 |q_i(Z,M)|\le\norm{(Y_Z)_i}\,\norm{(Y_M)_i},\qquad |q_i(M)|\le\norm{(Y_M)_i}^2 .
\end{equation}

\begin{lemma}\label{lem:commutator}
Let $f\in\R^n$, $F=\Diag(f)$, and $R=I-2P$ (an orthogonal matrix). Then:
\begin{enumerate}[label=(\alph*)]
\item $q(N_f)=\mathcal A\Gamma_f$ with $\Gamma_f=V^TFRFU$, and
$\fro{\Gamma_f}\le2\norm f_\infty\fro{N_f}$;
\item $T_f:=Y_{N_f}=FP+PF-2PFP$ satisfies
$[X,T_f]=RF[X,P]+[X,P]FR$ for diagonal $X$; hence
$\Lap(T_f)\preceq4\norm f_\infty^2L$, and
$x^T\Lap(T_{e_k})x\le2\sum_jP_{kj}^2(x_k-x_j)^2$.
\end{enumerate}
\end{lemma}

\begin{proof}
(a) For $x\in\R^n$ put $X=\Diag(x)$, $B_x=V^TXV$, $A_x=U^TXU$. Then
$\ip x{q(Z)}=\ip Z{B_xZ-ZA_x}$ for every $Z$, and, since $X$ and $F$ commute,
\[
 B_xN_f-N_fA_x=V^T[X(I-P)F-FPX]U=V^T[F(I-P)X-XPF]U=B_fN_x-N_xA_f .
\]
Hence $\ip x{q(N_f)}=\ip{N_f}{B_fN_x-N_xA_f}=\ip{B_fN_f-N_fA_f}{N_x}$. As
$\ip M{N_x}=\sum_ix_iv_i^TMu_i=\ip{\mathcal AM}x$ for every $M$, this says
$q(N_f)=\mathcal A(B_fN_f-N_fA_f)$, and
$B_fN_f-N_fA_f=V^TF(I-P)FU-V^TFPFU=\Gamma_f$. The norm bound follows from
$\opn{B_f},\opn{A_f}\le\norm f_\infty$.

(b) $Y_{N_f}=VV^TFUU^T+UU^TFVV^T=(I-P)FP+PF(I-P)$. Since $X$ commutes with $F$,
$[X,FP]=F[X,P]$, $[X,PF]=[X,P]F$ and $[X,PFP]=[X,P]FP+PF[X,P]$, which gives the
commutator formula. Then $\fro{[X,T_f]}\le2\norm f_\infty\fro{[X,P]}$, i.e.\ the
first bound (\cref{lem:energy}). For $f=e_k$ and $E_k=\Diag(e_k)$, $E_k[X,P]$ and $[X,P]E_k$ are
the $k$th row and column of $[X,P]$, each of squared norm
$\sum_jP_{kj}^2(x_k-x_j)^2$.
\end{proof}

\subsection{The second-order mean and the drift}

\begin{lemma}[Mean]\label{lem:mean}
$\E q(Z)=a\one-p-b_0-b_*$ with $b_*:=\mathcal Am_*$, where
\[
 b_0=\frac1nL(p^{-1}),\quad \norm{b_0}_\infty\le\frac{12\eps}n,\qquad
 m_*=\frac1n\sum_{\mu>0}\frac{r_\mu}\mu\,\Gamma_{f_y},\quad
 \fro{m_*}^2\le\frac{2a}\rho .
\]
\end{lemma}

\begin{proof}
For $G\sim N(0,I)$ on $\mathcal Z$, $\E\norm{G^Tv_i}^2=d(1-p_i)$ and
$\E\norm{Gu_i}^2=(n-d)p_i$, so the unfiltered covariance $I/n$ contributes
$a\one-p$. The covariance $\Omega=\sum_j\zeta_j\otimes\zeta_j$, $\zeta_j=\bar{\mathcal
A}^*e_j=v_ju_j^T/\sqrt{p_j}$, contributes
\[
 \frac1n\sum_jq_i(\zeta_j)=\frac1n\sum_j\frac{(I-P)_{ij}^2p_j-(1-p_j)P_{ij}^2}{p_j}
 =\frac1n\Bigl(1-\sum_j\frac{P_{ij}^2}{p_j}\Bigr)=\frac1n(Lp^{-1})_i ,
\]
using $\sum_j(I-P)_{ij}^2=1-p_i$ and $\sum_jP_{ij}^2=p_i$. Here
$(Lp^{-1})_i=\sum_jP_{ij}^2(p_i^{-1}-p_j^{-1})$ and
$|p_i^{-1}-p_j^{-1}|\le2\eps a/((1-\eps)a)^2\le8\eps/a$, so
$|(Lp^{-1})_i|\le12\eps$. Finally $R_\rho/n$ contributes
$\frac1n\sum_\mu r_\mu q(Z_y)=\frac1n\sum_\mu\frac{r_\mu}\mu q(N_{f_y})=\mathcal
Am_*$ by \cref{lem:commutator}(a), and with $\fro{N_{f_y}}=\sqrt\mu$,
\eqref{eq:finfty} and \cref{lem:filter}(c),
\[
 \fro{m_*}\le\frac2n\sqrt{\frac2a}\sum_{\mu>0}\frac{r_\mu}{\sqrt\mu}
 \le\frac2n\sqrt{\frac2a}\cdot\frac d{2\sqrt\rho}=\sqrt{\frac{2a}\rho}.\qedhere
\]
\end{proof}

\begin{remark}[Why a drift]\label{rem:drift}
The filter removes covariance from the directions $Z_y$, and with it their share
$t^2b_*=t^2\mathcal Am_*$ of the second-order drift towards $a\one$. This bias
need not be below the tolerance $\sigma/(12\Theta)$ used later: for instance, for
many small blocks joined through a common row it is of order $at^2$ at that row.
(Since $\Gamma_f=V^T[(ff^T)\circ(I-2P)]U$ is quadratic in $f$, and kernel modes
$Sy=0$ have $N_{f_y}=0$ and hence $\Gamma_{f_y}=0$, the spectral identity
$\sum_y\frac{1-\mu}{\rho+\mu}yy^T=(I-S)(\rho I+S)^{-1}$ shows that $m_*$ does not
depend on the choice of eigenbasis:
$m_*=\frac1nV^T\bigl[(D_p^{-1/2}(I-S)(\rho I+S)^{-1}D_p^{-1/2})\circ(I-2P)\bigr]U$.)
However, $t^2b_*$ is exactly the \emph{first-order} diagonal change produced
by the deterministic horizontal direction $\frac t2m_*$, whose norm is
$\frac t2\fro{m_*}\le\sqrt{a/2}/D$. Adding this direction to the noise cancels
the bias exactly, at the price of cross terms of size $O(\sigma/D)$.
\end{remark}

\subsection{The expected graph}

Since $\E Y_{ij}^2$ is linear in the covariance, $\E_\Sigma\Lap(Y)=\sum_k\lambda_k
\Lap(Y_{\psi_k})$ for every representation $\Sigma=\sum_k\lambda_k\psi_k\otimes\psi_k$
with real $\lambda_k$; we use it for $C_\rho/n=(I-\Omega-R_\rho)/n$ with
$\Omega=\sum_j\zeta_j\otimes\zeta_j$.

\begin{lemma}\label{lem:expected-graph}
$\displaystyle\E\Lap(Y)\succeq\frac a{4n}\mathsf K-\Bigl(\frac2n+\frac8d+\frac8\rho\Bigr)L$.
\end{lemma}

\begin{proof}
\emph{Unfiltered noise.} For $Z_0\sim N(0,I/n)$ and $i\ne j$,
$(Y_{Z_0})_{ij}=v_i^TZ_0u_j+v_j^TZ_0u_i$ has variance
$n^{-1}\fro{v_iu_j^T+v_ju_i^T}^2=n^{-1}(p_i+p_j-2p_ip_j-2P_{ij}^2)$. As
$p\le\frac34$, $p_i+p_j-2p_ip_j\ge\frac14(p_i+p_j)\ge\frac a4$, so
$\E\Lap(Y_{Z_0})\succeq\frac a{4n}\mathsf K-\frac2nL$.

\emph{Base loss.} $Y_{\zeta_k}=T_{e_k}/\sqrt{p_k}$, so by
\cref{lem:commutator}(b) and $p_k\ge a/2$,
$\sum_kx^T\Lap(Y_{\zeta_k})x\le\frac2a\cdot2\sum_{k,j}P_{kj}^2(x_k-x_j)^2
=\frac8a\,x^TLx$; the base part $\Omega/n$ therefore removes at most $\frac8dL$.

\emph{Residual loss.} $\Lap(Y_{Z_y})=\mu^{-1}\Lap(T_{f_y})\preceq\frac8{a\mu}L$ by
\cref{lem:commutator}(b) and \eqref{eq:finfty}, so $R_\rho/n$ removes at most
$\frac8{an}\sum_\mu\frac{r_\mu}\mu L\le\frac8\rho L$.
\end{proof}

\subsection{A good sample}

Let $\ell_i=\frac1n\sum_\mu r_\mu\norm{(Y_{Z_y})_i}^2$ be the variance that $R_\rho/n$
removes from row $i$. Since $\fro{Y_{Z_y}}^2=2$, $\sum_i\ell_i\le\frac2n\sum_\mu
r_\mu\le2a$. Fix $\delta_0=1/6400$ and let
\[
 \mathfrak B=\{i:\ell_i>\delta_0a\},\qquad b=|\mathfrak B|\le2/\delta_0=:b_{\max},
\]
a deterministic set depending only on $U$ and $\rho$. Write
$\mathcal C(Y)=\sum_{i<j}2P_{ij}Y_{ij}(e_i-e_j)(e_i-e_j)^T$, so that
$\Lap(P+tY)=L+t\mathcal C(Y)+t^2\Lap(Y)$.

\begin{lemma}\label{lem:sample}
There is an absolute $h\in(0,1)$ and, for each $\eta'\in(0,\frac1{32}]$, a
constant $A_{\eta'}$ such that the following holds for all $\rho>0$ and
$0<t\le1$ whenever $d\ge A_{\eta'}\log(2n)$. With positive probability:
\begin{enumerate}[label=\textup{(S\arabic*)}]
\item\label{S1} $\opn Z\le16$, $\max_i\norm{Z^Tv_i}^2\le3a$ and
$\max_i\norm{Y_i}^2\le800a$;
\item\label{S2} $\max_i|(\mathcal AZ)_i|\le3\sqrt{\rho a\log(2n)/n}$ and
$\max_i|q_i(Z)-\E q_i(Z)|\le\eta'a$;
\item\label{S3} every $i\notin\mathfrak B$ has at least $0.95n$ indices
$j\ne i$ with $|P_{ij}+tY_{ij}|\ge h\,t\sqrt{a/n}$;
\item\label{S4} for every $x$ supported on $\mathfrak B$,
\[
 |t\,x^T\mathcal C(Y)x|\le\eta'\bigl(x^TLx+at^2\norm x^2\bigr),\qquad
 |x^T(\Lap(Y)-\E\Lap(Y))x|\le\eta'a\norm x^2 .
\]
\end{enumerate}
\end{lemma}

\begin{proof}
Since $C_\rho\preceq I$, every linear functional $\ip Z\Psi$ has variance at
most $\fro\Psi^2/n$; we use the tools of \cref{app:gauss}.

\ref{S1}: The net bound (\cref{lem:gauss}(c)) with $\sigma^2=1/n$ gives
$\Prob(\opn Z>16)\le2\cdot5^ne^{-8n}$. The vector $Z^Tv_i$ has covariance
$\preceq n^{-1}I_d$ and trace $\le a$; \cref{lem:gauss}(a) with $u=2a$ (so
$\min\{u^2/\fro\cdot^2,u/\opn\cdot\}\ge2d$) and a union bound give
$\max\norm{Z^Tv_i}^2\le3a$ off probability $2ne^{-cd}$. Then
$\norm{Y_i}\le\norm{Z^Tv_i}+\norm{Zu_i}\le\sqrt{3a}+16\sqrt{3a/2}$.

\ref{S2}: $(\mathcal AZ)_i$ is Gaussian with variance $\le\rho p_i/n\le3\rho
a/(2n)$ (\cref{lem:filter}(b)); a union bound gives the first claim off
probability $2n(2n)^{-3}$. The form $q_i(Z)=\ip Z{\mathsf M_iZ}$,
$\mathsf M_iZ=v_iv_i^TZ-Zu_iu_i^T$, has $\opn{\mathsf M_i}\le1$ and
$\fro{\mathsf M_i}^2\le2(d+np_i^2)\le7d$; after the covariance factor
$(C_\rho/n)^{1/2}$ these become $\le1/n$ and $\le\sqrt{7d}/n$, and
\cref{lem:gauss}(a) with $u=\eta'a$ gives failure probability $2n\exp(-c\eta'^2d)$.

\ref{S3}: \emph{Variance bookkeeping.} By the proof of
\cref{lem:expected-graph}, for $i\ne k$ with $P_{ik}^2\le a/16$ the unfiltered
variance of $Y_{ik}$ is at least $a/(8n)$; at most $24$ indices $k$ have
$P_{ik}^2>a/16$. Row $i$ of $T_{e_j}$ is
$\delta_{ij}e_j^TP+P_{ij}e_j^T-2P_{ij}e_j^TP$, of squared norm at most
$3(\delta_{ij}p_j+P_{ij}^2+4P_{ij}^2p_j)$, so the base part removes from row $i$
total variance $\frac1n\sum_j\norm{e_i^TT_{e_j}}^2/p_j\le\frac6{an}\cdot5p_i\le
\frac{45}n$, which exceeds $\frac a{32n}$ at no more than $1440n/d$ positions.
For $i\notin\mathfrak B$ the residual part removes at most $\delta_0a$ from row
$i$, hence more than $\frac a{32n}$ at no more than $32\delta_0n$ positions. So for
$d\ge10^6$ (hence $n\ge2\cdot10^6$, and the excluded indices number at most
$25+0.0065n\le0.01n$) every $i\notin\mathfrak B$ has at least $0.99n$ indices $k\ne i$
with $\Var(Y_{ik})\ge\frac a{16n}$.

\emph{Small ball.} Let $\phi:\R\to[0,1]$ be $C^1$, equal to $1$ on $[-h,h]$,
vanishing outside $[-2h,2h]$, with $|\phi'|\le2/h$, and for $i\notin\mathfrak B$
let $F_i=\frac1n\sum_{k\ne i}\phi\bigl((P_{ik}+tY_{ik})/(t\sqrt{a/n})\bigr)$. By
the density bound for a Gaussian of variance $\ge\frac{at^2}{16n}$,
$\E F_i\le0.01+\frac{4h}{\sqrt{2\pi/16}}\le0.01+7h$. As a function of the row
$(Y_{ik})_k$, $F_i$ has gradient norm at most
$\frac1n\cdot\frac2h\cdot\sqrt{n/a}\cdot\sqrt n=2/(h\sqrt a)$; the row depends on
$Z$ through a linear map of norm $\le2$, and $Z=(C_\rho/n)^{1/2}g$ with $g$
standard. So $F_i$ is $4/(h\sqrt d)$-Lipschitz in $g$, and \cref{lem:gauss}(b)
gives $\Prob(F_i>\E F_i+0.01)\le2\exp(-c\,h^2d)$. Take $h=1/500$, so that
$\E F_i+0.01\le0.04$, and a union bound over $i$. On the complement, at most
$0.04n$ indices $k\ne i$ have $|P_{ik}+tY_{ik}|<ht\sqrt{a/n}$, so at least
$n-1-0.04n\ge0.95n$ have the reverse inequality.

\ref{S4}: Only second moments are needed, because $\mathfrak B$ is deterministic
and bounded. Let $\Pi_{\mathfrak B}$ be the coordinate projection onto
$\R^{\mathfrak B}$, $J_{\mathfrak B}=L_{\mathfrak B\mathfrak B}+at^2I_{\mathfrak B}$
and $\omega_{ij}=J_{\mathfrak B}^{-1/2}\Pi_{\mathfrak B}(e_i-e_j)$, so
$\norm{\omega_{ij}}^2\le2/(at^2)$ and
$\sum_{i<j}P_{ij}^2\omega_{ij}\omega_{ij}^T=J_{\mathfrak B}^{-1/2}L_{\mathfrak
B\mathfrak B}J_{\mathfrak B}^{-1/2}\preceq I$. For $x$ supported on $\mathfrak
B$, $tx^T\mathcal C(Y)x=\ip{J_{\mathfrak B}^{1/2}x}{\Xi J_{\mathfrak B}^{1/2}x}$
with $\Xi=\sum_{i<j}2tP_{ij}Y_{ij}\,\omega_{ij}\omega_{ij}^T$, and
$x^TJ_{\mathfrak B}x=x^TLx+at^2\norm x^2$. The coefficients $(Y_{ij})_{i<j}$
have covariance $\preceq\frac2nI$ (the map $Z\mapsto Y_Z$ is $\sqrt2$ times an
isometry), so
\[
 \E\fro\Xi^2\le\frac{2t^2}n\sum_{i<j}4P_{ij}^2\norm{\omega_{ij}}^4
 \le\frac{16}{an}\tr\bigl(J_{\mathfrak B}^{-1/2}L_{\mathfrak B\mathfrak B}
 J_{\mathfrak B}^{-1/2}\bigr)\le\frac{16b}d ,
\]
and Markov's inequality gives the first half off probability $16b/(d\eta'^2)$.
The principal $\mathfrak B$-block of $\Lap(Y)-\E\Lap(Y)$ has diagonal entries
$\sum_{j\ne i}(Y_{ij}^2-\E Y_{ij}^2)$, a centred Gaussian quadratic form of
variance at most $2\cdot\frac2n\cdot\E\norm{Y_i}^2\le\frac{12a}n$
(\cref{lem:gauss}(a), with $\E\norm{Y_i}^2\le a(1-p_i)+(1-a)p_i\le3a$), and
off-diagonal entries $-(Y_{ik}^2-\E Y_{ik}^2)$ of variance at most $8/n^2$. Its
expected squared Frobenius norm is at most $20b^2a/n=20b^2a^2/d$, and Markov
gives the second half off probability $20b^2/(d\eta'^2)$.

Finally choose $A_{\eta'}$ so large that $d\ge A_{\eta'}\log(2n)$ implies
$d\ge10^6$ and all failure probabilities above are at most $\frac1{10}$; their
sum is less than $1$.
\end{proof}

\subsection{The drifted retraction}

Fix $D\ge12$ and $t\in(0,1]$, set $\rho=D^2t^2$, take a sample $Z$ as in
\cref{lem:sample}, and let $m_*$ be as in \cref{lem:mean}. Define
\[
 W=Z+\tfrac t2m_*,\qquad G=I+t^2W^TW,\qquad U_t=(U+tVW)G^{-1/2},\qquad
 P_t=U_tU_t^T .
\]
Since $U^TV=0$, $(U+tVW)^T(U+tVW)=G$ and $U_t$ is Parseval.

\begin{lemma}\label{lem:retraction}
There is an absolute $C_E$ (depending only on the bounds in \ref{S1} and on
$D\ge12$) such that:
\begin{enumerate}[label=(\alph*)]
\item $\fro{U_t-U}^2\le C_Et^2d$;
\item every row of $\mathsf E=P_t-P-tY$ satisfies
$\sum_j\mathsf E_{ij}^2\le at^2(D^{-2}+C_Et^2)$;
\item with $\norm{\mathsf r}_\infty\le C_Eat^3$,
\[
 \diag P_t-a\one=(1-t^2)(p-a\one)+2t\mathcal AZ-t^2b_0+t^2\bigl(q(Z)-\E q(Z)\bigr)
 +t^3q(Z,m_*)+\tfrac{t^4}4q(m_*)+\mathsf r ;
\]
\item $\norm{\diag P_t-a\one}_\infty\le at^2\bigl(\frac{\eps}{t^2}
+6D\sqrt{\log(2n)/d}+\frac{12\eps}d+\eta'+\frac{40}D+\frac1{2D^2}+C_Et\bigr)$.
\end{enumerate}
\end{lemma}

\begin{proof}
By \cref{lem:mean}, $\frac t2\opn{m_*}\le\frac t2\fro{m_*}\le\sqrt{a/2}/D\le\frac1{24}$,
so with \ref{S1}: $\opn W\le17$, $\fro W\le17\sqrt d$,
$\norm{W^Tv_i}\le2\sqrt a$, $\norm{Wu_i}\le17\sqrt{3a/2}$, and
$\norm{(Y_{m_*})_i}^2\le\opn{m_*}^2\le2a/\rho$. Hence $\opn{G^{-1}-I}$ and
$\opn{G^{-1/2}-I}$ are at most $289t^2$, $\opn{U+tVW}\le18$ and
$\norm{u_i+tW^Tv_i}\le4\sqrt a$.

(a) $\fro{U_t-U}\le\fro{U+tVW}\opn{G^{-1/2}-I}+t\fro W\le Ct\sqrt d$.

(b) $\mathsf E=\frac{t^2}2Y_{m_*}+\mathsf E'$ with
$\mathsf E'=P_t-P-tY_W=(U+tVW)(G^{-1}-I)(U+tVW)^T+t^2VWW^TV^T$, whose rows have
norm at most $Ct^2\sqrt a$. The rows of $\frac{t^2}2Y_{m_*}$ have squared norm at
most $\frac{t^4}4\cdot\frac{2a}\rho=\frac{at^2}{2D^2}$. Use $(\alpha+\beta)^2\le
2\alpha^2+2\beta^2$.

(c) Let $\xi_i=u_i+tW^Tv_i$, the $i$th row of $U+tVW$. Inserting
$G^{-1}=I-t^2W^TW+t^4(W^TW)^2G^{-1}$ into $(P_t)_{ii}=\xi_i^TG^{-1}\xi_i$ and using
$u_i^TW^Tv_i=(\mathcal AW)_i$,
\[
 (P_t)_{ii}=p_i+2t(\mathcal AW)_i+t^2q_i(W)-2t^3\ip{Wu_i}{WW^Tv_i}
 -t^4\norm{WW^Tv_i}^2+t^4\xi_i^T(W^TW)^2G^{-1}\xi_i ,
\]
and the last three terms are $O(at^3)$. Now $2t\mathcal AW=2t\mathcal AZ+t^2\mathcal
Am_*$, $q(W)=q(Z)+tq(Z,m_*)+\frac{t^2}4q(m_*)$, and
$\E q(Z)=a\one-p-b_0-\mathcal Am_*$ by \cref{lem:mean}; the terms $\pm t^2\mathcal
Am_*$ cancel.

(d) In (c): $(1-t^2)\norm{p-a\one}_\infty\le\eps a$;
$2t\norm{\mathcal AZ}_\infty\le6tD t\sqrt{a\log(2n)/n}$ by \ref{S2}, which is
$6Dat^2\sqrt{\log(2n)/d}$; $t^2\norm{b_0}_\infty\le12\eps t^2/n$; the fluctuation is
at most $\eta'at^2$; by~\eqref{eq:q-cs} and \ref{S1},
$t^3|q_i(Z,m_*)|\le t^3\sqrt{800a}\sqrt{2a/\rho}=40at^2/D$ and
$\frac{t^4}4|q_i(m_*)|\le\frac{t^4}4\cdot\frac{2a}\rho=\frac{at^2}{2D^2}$.
\end{proof}

\begin{lemma}[Graph of the seed]\label{lem:seed-graph}
There are absolute $t_*\in(0,1]$ and $C_{\rm core}$ such that the following
holds. Let $\eta'\le\frac1{32}$, $d\ge A_{\eta'}\log(2n)$, $D\ge13/h$,
$b\le n/10$ and $t\le t_*$, and let $Z$ satisfy \ref{S1}--\ref{S4}. Then
$H(L_{P_t})\le C_{\rm core}/(at^2)$.
\end{lemma}

\begin{proof}
\emph{Core.} In a row $i\notin\mathfrak B$, by \cref{lem:retraction}(b), at most
$at^2(D^{-2}+C_Et^2)\big/\bigl(\frac{h^2}4\frac{at^2}n\bigr)\le0.05n$ indices have
$|\mathsf E_{ij}|\ge\frac h2t\sqrt{a/n}$, once $D\ge13/h$ and
$t^2\le h^2/(160C_E)$. With \ref{S3}, every $i\notin\mathfrak B$ has at least
$0.9n$ indices $j\ne i$ with $(P_t)_{ij}^2\ge\frac{h^2}4\frac{at^2}n$.

\emph{Block.} For $x$ supported on $\mathfrak B$, \cref{lem:expected-graph},
$x^T\mathsf Kx\ge(n-b)\norm x^2$ and \ref{S4} give
\[
 x^T\Lap(P+tY)x\ge\Bigl(1-\eta'-t^2\bigl(\tfrac2n+\tfrac8d\bigr)-\tfrac8{D^2}\Bigr)x^TLx
 +at^2\Bigl(\tfrac14\bigl(1-\tfrac bn\bigr)-2\eta'\Bigr)\norm x^2
 \ge\frac{at^2}8\norm x^2 ,
\]
since the first coefficient is positive and $\frac14\cdot\frac9{10}-\frac1{16}\ge\frac18$.
By \eqref{eq:lap-basic} and \cref{lem:retraction}(b),
$L_{P_t}=\Lap(P+tY+\mathsf E)\succeq\frac12\Lap(P+tY)-2at^2(D^{-2}+C_Et^2)I$,
which is at least $\frac{at^2}{32}$ on vectors supported on $\mathfrak B$ once
$t^2\le1/(128C_E)$.

\emph{Barrier.} \Cref{lem:core} with $\mathcal B=\mathfrak B$, part (i) with
$\vartheta=0.4$, $\gamma=h^2at^2/4$, and part (ii) with $\lambda=at^2/32$, gives
$H(L_{P_t})\le\frac{1+b}{0.1h^2at^2}+\frac{32\sqrt b}{at^2}$. Put
$C_{\rm core}=10(1+b_{\max})/h^2+32\sqrt{b_{\max}}$.
\end{proof}

\subsection{Proof of \texorpdfstring{\cref{thm:moderate}}{Theorem}}

\emph{Constants.} The numbers $\delta_0,b_{\max},h,C_E,t_*,C_{\rm core}$ are
absolute and do not depend on $D$ (for $D\ge12$). Put
$\Theta=\max\{2,C_{\rm core},C_E\}$ and choose, in this order: $D\ge\max\{13/h,
500\Theta\}$, so that $\frac{40}D+\frac1{2D^2}\le\frac1{12\Theta}$; $K_0=12\Theta$
and $\eta'=\min\{\frac1{32},\frac1{12\Theta}\}$; $A\ge A_{\eta'}$ so large that
$d\ge A\log(2n)$ implies $6D\sqrt{\log(2n)/d}\le\frac1{12\Theta}$ and
$b_{\max}\le n/10$; and $\eps_0\le\frac12$ so that $t=\sqrt{K_0\eps}$ satisfies
$t\le t_*$, $C_Et\le\frac1{12\Theta}$ and $12\eps\le\frac1{12\Theta}$ for
$\eps\le\eps_0$. None of these choices depends on $n,d,\eps$ or $U$.

\emph{Seed.} With $t^2=K_0\eps$, \cref{lem:retraction}(d) bounds the diagonal
error of $P_t$ by $at^2\cdot\frac6{12\Theta}=\frac{at^2}{2\Theta}$;
\cref{lem:seed-graph} bounds the barrier constant by $\Theta/(at^2)$; and
\cref{lem:retraction}(a) bounds $\fro{U_t-U}^2$ by $\Theta t^2d$. So $U_t$ is a
$(\Theta,t)$-seed for $U$, and \cref{cor:seed} gives an ENP frame $\widehat U$
with $\fro{\widehat U-U}^2\le3\Theta t^2d=36\Theta^2\eps d$.
\qed

\begin{remark}
The drift keeps the spectral input local. The bias $t^2b_*$ could also be
removed after the perturbation, for instance by a second scaling; since
$\ip x{b_*}^2\le\frac{2a}\rho x^TLx$ this is possible, but it requires two-sided
control of $L_{P_t}$ on all of $\one^\perp$. With the drift, the only spectral
input is the gap on the bounded deterministic block $\mathfrak B$, which needs
second moments only.
\end{remark}

\section{Assembly}\label{sec:assembly}

We use the two seeds through the following equal-norm formulation.

\begin{proposition}\label{prop:equalrow}
There is an absolute constant $C_0$ such that every $X\in\R^{n\times d}$ with
$1\le d$, $2d\le n$ and $\norm{x_i}^2=a$ for all $i$ admits an ENP frame $W$
with
\[
 \fro{X-W}^2\le C_0\,d\,\opn{X^TX-I}.
\]
\end{proposition}

\begin{proof}
Let $\eta=\opn{X^TX-I}$; if $\eta=0$ take $W=X$. Let $A,\eps_0,C_m$ be the
constants of \cref{thm:moderate} and $B,c_*,C_*$ those of \cref{thm:manyrow},
and fix an integer $d_0\ge2$ with $A\log(2Bd^2)\le d$ for all $d\ge d_0$.

\emph{Case $d<d_0$.} \Cref{lem:HM} gives $\fro{X-W}^2\le\eta d(d-1)\le d_0\eta d$.

\emph{Case $d\ge d_0$, $n\ge Bd^2$.} If $\eta\le c_*$ use \cref{thm:manyrow};
otherwise \cref{cor:exist} gives cost $4d\le(4/c_*)\eta d$.

\emph{Case $d\ge d_0$, $n<Bd^2$.} Then $A\log(2n)\le d$. If $\eta\le\eps_0/4$, let
$U=X(X^TX)^{-1/2}$. The singular values of $X$ lie in
$[\sqrt{1-\eta},\sqrt{1+\eta}]$, so $\fro{X-U}^2\le d\eta^2$, and
$\norm{u_i}^2=x_i^T(X^TX)^{-1}x_i\in[a/(1+\eta),a/(1-\eta)]$, so
$|\norm{u_i}^2-a|\le2\eta a$. \Cref{thm:moderate} with $\eps=2\eta$ gives an ENP
$W$ with $\fro{U-W}^2\le2C_m\eta d$, and
$\fro{X-W}^2\le2d\eta^2+4C_m\eta d$. If $\eta>\eps_0/4$ use the trivial cost
$4d\le(16/\eps_0)\eta d$. Altogether
$C_0=\max\{d_0,C_*,4/c_*,2+4C_m,16/\eps_0\}$ works.
\end{proof}

\begin{proof}[Proof of \cref{thm:projection}]
Replacing $(P,Q)$ by $(I-P,I-Q)$ preserves $\beta$, ranks of complements and
distances, so we may assume $d\le n/2$; the case $d=0$ is trivial. Write
$P=UU^T$ with $U^TU=I$, $p_i=\norm{u_i}^2$ and $\delta=\beta/a$. If
$\delta\ge\frac12$, any ENP $W$ (\cref{cor:exist}) gives
$\fro{P-WW^T}^2\le2d\le4n\beta$. Otherwise let $x_i=\sqrt a\,u_i/\norm{u_i}$. Then
$\fro{X-U}^2=\sum_i(\sqrt{p_i}-\sqrt a)^2\le\sum_i(p_i-a)^2/a\le\delta^2d$ and
$X^TX=U^T\Diag(a/p)U$, so $\opn{X^TX-I}\le\max_i|a/p_i-1|\le\delta/(1-\delta)\le
2\delta$. \Cref{prop:equalrow} gives an ENP $W$ with
$\fro{U-W}^2\le2\delta^2d+4C_0\delta d\le(1+4C_0)n\beta$, using $\delta d=n\beta$.
Take $Q=WW^T$; by \cref{lem:align}, $\fro{P-Q}^2\le2(1+4C_0)n\beta$. So
\cref{thm:projection} holds with $C_P=\max\{4,2+8C_0\}$.
\end{proof}

\begin{proof}[Proof of \cref{thm:main}]
If $\eps\ge\frac12$, an ENP $W$ gives $\fro{X-W}^2\le2\fro X^2+2d\le6d\le12\eps
d$, since $\fro X^2=\tr X^TX\le(1+\eps)d$. Let $\eps<\frac12$ and
$U=X(X^TX)^{-1/2}$, $P=UU^T$. As above $\fro{X-U}^2\le\eps^2d$ and
$P_{ii}=x_i^T(X^TX)^{-1}x_i\in[\frac{1-\eps}{1+\eps}a,\frac{1+\eps}{1-\eps}a]$, so
$\beta=\max_i|P_{ii}-a|\le4\eps a$. \Cref{thm:projection} gives $Q$ with
$\fro{P-Q}^2\le4C_P\eps d$, and by \cref{lem:align} there is an isometry $W$
spanning $Q$---hence ENP---with $\fro{U-W}^2\le\fro{P-Q}^2$. Then
$\fro{X-W}^2\le2\eps^2d+8C_P\eps d$, so \cref{thm:main} holds with
$C=\max\{12,1+8C_P\}$.
\end{proof}

\appendix
\section{Gaussian tools}\label{app:gauss}

\begin{lemma}\label{lem:gauss}
The following hold with absolute constants.
\begin{enumerate}[label=(\alph*)]
\item \emph{(Quadratic forms.)} For $g$ standard Gaussian and $\mathsf A$
symmetric, $\Var(g^T\mathsf Ag)=2\fro{\mathsf A}^2$ and
\[
 \Prob\{|g^T\mathsf Ag-\tr\mathsf A|>u\}\le2\exp\Bigl[-\tfrac18\min\Bigl\{\frac{u^2}{\fro{\mathsf A}^2},\frac u{\opn{\mathsf A}}\Bigr\}\Bigr].
\]
For a Gaussian vector $\zeta=\Sigma^{1/2}g$, apply this to
$\Sigma^{1/2}\mathsf A\Sigma^{1/2}$, whose norms are at most $\opn\Sigma$ times those of
$\mathsf A$.
\item \emph{(Lipschitz functions.)} If $F$ is $C^1$ with $\norm{\nabla F}\le\ell$
and $g$ is standard Gaussian, then for $u\ge0$
$\Prob\{|F(g)-\E F(g)|>u\}\le2\exp(-2u^2/(\pi^2\ell^2))$.
\item \emph{(Nets.)} If $Z$ is an $r\times s$ centred Gaussian matrix with
$\Cov(\mathrm{vec}\,Z)\preceq\sigma^2I$, then for $u\ge0$
$\Prob\{\opn Z>u\}\le2\cdot5^{r+s}\exp(-u^2/(32\sigma^2))$.
\end{enumerate}
\end{lemma}

\begin{proof}
(a) Diagonalise $\mathsf A$ with eigenvalues $\lambda_j$. For $|s|\le1/(4\opn{\mathsf A})$,
$\log\E e^{s(g^T\mathsf Ag-\tr\mathsf A)}=\sum_j[-s\lambda_j-\frac12\log(1-2s\lambda_j)]
\le2s^2\fro{\mathsf A}^2$, since $-x-\frac12\log(1-2x)=\sum_{k\ge2}(2x)^k/(2k)\le2x^2$
for $|x|\le\frac14$. The exponential Markov inequality with
$s=\min\{u/(4\fro{\mathsf A}^2),1/(4\opn{\mathsf A})\}$, applied to both signs,
gives the tail. The variance follows from $\E g_j^4=3$.

(b) Let $g'$ be an independent copy, $g_\theta=g\cos\theta+g'\sin\theta$ and
$\dot g_\theta=-g\sin\theta+g'\cos\theta$; for each $\theta$ the pair
$(g_\theta,\dot g_\theta)$ is standard Gaussian. By the fundamental theorem of
calculus and Jensen on $[0,\pi/2]$,
\[
 \E e^{s(F(g')-F(g))}\le\frac2\pi\int_0^{\pi/2}\E\exp\Bigl[\frac{\pi s}2\ip{\nabla
 F(g_\theta)}{\dot g_\theta}\Bigr]d\theta\le e^{\pi^2s^2\ell^2/8}.
\]
Jensen in $g'$ gives $\E e^{s(F(g)-\E F)}\le e^{\pi^2s^2\ell^2/8}$; optimise in $s$.
All expectations are finite since $|F(x)|\le|F(0)|+\ell\norm x$, and Fubini
applies to the $\theta$-integral.

(c) Let $\mathcal N_1,\mathcal N_2$ be $\frac12$-nets of the unit spheres of $\R^r$
and $\R^s$ with at most $5^r$, $5^s$ points. For a vector $w$, a net point $y'$
within $\frac12$ of $w/\norm w$ has $\ip w{y'}\ge\frac12\norm w$; applying this
twice gives $\opn Z\le4\max_{x\in\mathcal N_1,y\in\mathcal N_2}|x^TZy|$. Each
$x^TZy$ is Gaussian with variance $\le\sigma^2$; a union bound finishes.
\end{proof}

\paragraph{Covariance domination.} If a Gaussian vector $\xi$ has covariance
$\Gamma\preceq\gamma I$ and $(M_\alpha)$ are matrices, then
$\E\fro{\sum_\alpha\xi_\alpha M_\alpha}^2=\tr(\Gamma\mathsf G)\le\gamma\sum_\alpha
\fro{M_\alpha}^2$, where $\mathsf G\succeq0$ is the Gram matrix of the $M_\alpha$.
Correlated Gaussians enter only as $\Sigma^{1/2}g$: in quadratic forms through the
last sentence of (a) (\cref{lem:Q}, \cref{lem:sample}), in Lipschitz functions
through (b), in operator norms through (c), and in Frobenius norms of linear
combinations through this identity (\cref{lem:sample}\ref{S4}).

\section{Formalisation}\label{app:formal}

\Cref{thm:main} is formalised in Lean~4 with Mathlib; the development
accompanies this paper (directory \texttt{lean/}). The formal statement is
\begin{center}
\begin{minipage}{0.92\textwidth}
\small
\begin{verbatim}
def SharpPaulsenBound : Prop :=
  ∃ C : ℝ, 0 < C ∧
    ∀ (n d : ℕ), 0 < d → d ≤ n →
    ∀ (ε : ℝ) (U : Frame n d), 0 < ε → ε < 1 →
      IsNearlyEqualNormParseval ε U →
      ∃ W : Frame n d, IsEqualNormParseval W ∧
        sqDistance U W ≤ C * ε * (d : ℝ)
\end{verbatim}
\end{minipage}
\end{center}
where \texttt{Frame n d} is the type of real $n\times d$ matrices, Parsevalness
is $U^TU=1$, the near-Parseval condition is stated through the quadratic form
$x\mapsto\norm{Ux}^2$, the row-norm conditions refer to $d/n$ as a real
quotient, and \texttt{sqDistance} is the unnormalised squared Frobenius
distance. The projection form (\cref{thm:projection}) is derived for arbitrary
real symmetric idempotent matrices of rank $d$. Both statements are also
restated in plain Mathlib vocabulary (positive semidefiniteness, dot products,
traces, \texttt{Matrix.rank}), and the restatements are derived from them.

FORMAL-STATUS-PLACEHOLDER


\section*{Contributions and acknowledgements}
Hugo Eberhard and Kaarel H\"anni proposed constrained random perturbations at
squared cost $O(\eps d)$ as a starting strategy. Their calculations identified
the squared-Gram Laplacian governing infinitesimal scaling and the interaction
scale $\eps d/n^2$ that a perturbation should create between nearly separate
components. The subsequent control of the quadratic bias and the quantitative
balancing argument were developed in the course of completing the proof.

Astra, working through OpenAI's Codex system with collaborating agents,
developed the initial complete proof and its Lean formalisation. This work
included covariance filtering, bias correction, both perturbation regimes,
static balancing controlled by half-set hitting times, and the global
assembly. Claude introduced the cancelling drift in the moderate-row seed
and the centred remainder estimate that removes the dimension-dependent error
restriction in the many-row seed, thereby eliminating the random partition.
Claude also recast the balancing and dense-core arguments using explicit
barriers, prepared the streamlined manuscript, and formalised that version
with its stated constants. In subsequent revisions, Astra simplified the
scaling-distance estimate by a trace identity, replaced the covariance by an
ordinary resolvent, shortened the optimality argument, and formalised these
changes.

The proof builds on the perturbation-and-scaling approach of Kwok, Lau, Lee
and Ramachandran~\cite{KLLR}, the projection formulation of Cahill and
Casazza~\cite{CC}, and the quadratic bound of Hamilton and Moitra~\cite{HM}.
Ramachandran's dissertation and the work of Lau and
Ramachandran~\cite{RamThesis,LR} also informed the investigation. Neither the
dissertation's Paulsen results nor the announced linear bound is assumed in
this proof. These acknowledgements describe
the development of this argument and its intellectual influences; no claim of
historical priority is intended.

\begin{thebibliography}{9}
\bibitem{CC}
J.~Cahill and P.~G.~Casazza.
\emph{The Paulsen problem in operator theory}.
Operators and Matrices 7 (2013).
\href{https://arxiv.org/abs/1102.2344}{arXiv:1102.2344}.
\bibitem{HM}
L.~Hamilton and A.~Moitra.
\emph{The Paulsen problem made simple}.
ITCS 2019; Israel J.\ Math.\ (2021).
\href{https://arxiv.org/abs/1809.04726}{arXiv:1809.04726}.
\bibitem{KLLR}
T.~C.~Kwok, L.~C.~Lau, Y.~T.~Lee and A.~Ramachandran.
\emph{The Paulsen problem, continuous operator scaling, and smoothed analysis}.
STOC 2018.
\href{https://arxiv.org/abs/1710.02587}{arXiv:1710.02587}.
\bibitem{LR}
L.~C.~Lau and A.~Ramachandran.
\emph{Optimal bounds for Tyler's M-estimator for elliptical distributions},
Section~7.
\href{https://arxiv.org/abs/2510.13751}{arXiv:2510.13751}, 2025.
\bibitem{RamThesis}
A.~Ramachandran.
\emph{Geodesic convex analysis of group scaling for the Paulsen problem
and the tensor normal model}.
Ph.D. thesis, University of Waterloo, 2021.
\href{https://uwspace.uwaterloo.ca/handle/10012/17710}{University repository}.
\end{thebibliography}
\end{document}
